Fact verification, calibration, uncertainty quantification, educational AI, selective prediction, temperature scaling, ensemble methods, reproducibility

\section{Introduction}\label{introduction}

fact verification systems face a fundamental trade-off: achieving high accuracy while maintaining transparency about uncertainty. Students and instructors need not only to know \emph{if} a claim is supported or refuted, but also \emph{how confident} the system is in its judgment. Overconfident predictions can mislead learners, while overly cautious systems that abstain on most claims provide little value.

Modern large language models (LLMs) have demonstrated impressive capabilities in educational contexts, yet their tendency toward overconfident predictions poses risks in fact-checking applications~. A system claiming 95\% confidence when its true accuracy is 75\% can undermine trust and lead to pedagogical harm. This problem is particularly acute in educational settings where: students may lack domain expertise to question confident but incorrect predictions; instructors need reliable confidence scores to decide when to intervene; real-time feedback loops require sub-second latency; and deployment environments often lack access to commercial APIs.

To address these challenges, this paper presents CalibraTeach, a calibrated selective prediction system specifically designed for educational fact verification. We make five key contributions spanning methodological innovation, system engineering, resource release, empirical insights, and responsible deployment:

\begin{enumerate}
\def\labelenumi{\arabic{enumi}.}
\item
  \textbf{Methodological Contribution---Calibration Parity Protocol:} A systematic comparative evaluation methodology ensuring all baseline comparisons use identical calibration procedures (temperature scaling on the same validation set), isolating architectural improvements from calibration effects and enabling apples-to-apples fairness in evaluation.
\item
  \textbf{Systems Contribution---Real-Time Multi-Component Pipeline:} A 6-signal ensemble architecture achieving mean end-to-end latency of 67.68 ms (14.78 claims/sec throughput on NVIDIA RTX 4090), combining semantic relevance, entailment strength, evidence diversity, agreement, margin, and source authority with learned weights, complete with stage-by-stage latency breakdown and deployment-ready implementation.
\item
  \textbf{Resource Contribution---CSClaimBench and Reproducibility Artifacts:} A 1,045-claim expert-annotated dataset spanning five computer science subdomains (\(`\kappa=0.89`\)), curated evidence corpus (12,500 documents), and complete artifact generation pipeline producing 31 analysis files with deterministic evaluation validation and comprehensive confidence intervals (2000-sample stratified bootstrap, BCa method).
\item
  \textbf{Empirical Insight---Calibration as Control Signal for Abstention:} Demonstrate that well-calibrated confidence enables selective prediction to achieve 74\% automated coverage at 90\% selective accuracy, deferring uncertain cases to instructors with 92\% instructor agreement on abstention recommendations (\(`n=25`\) pilot), establishing calibration quality as the foundation for human-AI collaboration in education.
\item
  \textbf{Responsible Deployment Framework---Honest Limitations and RCT Imperative:} Explicit disclosure of 7 limitations including domain specificity (computer science claims only), sample size constraints (260-claim test set), and the critical caveat that pedagogical effectiveness claims require randomized controlled trial validation---technical feasibility does not imply learning benefits.
\end{enumerate}

\emph{Core insight}: In educational fact verification, calibrated confidence is not merely a metric to optimize---it serves as the control signal for abstention and instructor oversight. Well-calibrated predictions enable systems to defer uncertain cases reliably, transforming verification from a fully automated task into a principled human-AI collaborative workflow.

\section{Related Work}\label{related-work}

\textbf{Fact Verification Systems.}

The FEVER dataset~ established a benchmark for fact extraction and verification with 185,000 claims annotated against Wikipedia. While influential, FEVER focuses on general-domain claims and does not address educational contexts or calibration. SciFact~ extended fact verification to scientific claims, demonstrating domain-specific challenges in biomedical literature. ClaimBuster~ pioneered claim detection in political speeches but did not provide confidence calibration.

Recent work has explored neural architectures for fact verification, including BERT-based models~ and retrieval-augmented generation (RAG)~. Large language models (LLMs) have shown promise for zero-shot fact-checking~, but often exhibit poor calibration without explicit uncertainty quantification. Multi-hop reasoning systems~ address complex claims requiring multiple evidence sources, but computational costs limit real-time deployment.

\textbf{Calibration in Neural Networks.}

Temperature scaling~ remains the most effective post-hoc calibration method for deep networks, learning a single scalar parameter to rescale logits. Platt scaling~ applies logistic regression to calibrate binary classifiers. More complex methods like isotonic regression~ and Bayesian Binning into Quantiles~ offer flexibility but risk overfitting on small validation sets.

Ensemble diversity improves calibration~, though deep ensembles incur high computational costs. Recent work on conformal prediction~ provides distribution-free uncertainty quantification with theoretical coverage guarantees, but requires held-out calibration data.

For fact verification, calibration has received limited attention. Most systems report only accuracy and F1 scores~, omitting expected calibration error (ECE) or reliability diagrams. CalibraTeach addresses this gap by treating calibration as a first-class design objective.

\textbf{Why Temperature Scaling + Selective Prediction instead of Conformal Prediction?}

Recent work on conformal prediction~ offers distribution-free uncertainty quantification with theoretical coverage guarantees. CalibraTeach employs temperature scaling (post-hoc calibration) coupled with selective prediction (threshold-based abstention) rather than conformal prediction for the following reasons: \emph{(1)} Our system design requires a single scalar confidence threshold \(`\tau`\) for abstention decisions and instructor UI, which aligns naturally with temperature-scaled logits; conformal prediction would require representing prediction sets, adding UI complexity. \emph{(2)} We did not implement conformal prediction in this work; our approach uses empirical validation-set calibration, which is well-established for small-to-medium validation sets (\(`n=261`\)). \emph{(3)} We acknowledge that conformal methods offer theoretical advantages (distribution-free coverage guarantees) worth investigating on larger validation sets and alternative deployment scenarios. Accordingly, stochastic robustness evaluation and conformal prediction evaluation are deferred to future work (Section~8.1).

\textbf{Selective Prediction and Abstention.}

Selective prediction (also called prediction with rejection)~ allows models to abstain on uncertain inputs. Recent work~ demonstrates that selective classifiers can achieve high precision on retained predictions by trading off coverage. The coverage-accuracy trade-off is formalized through risk-coverage curves~, enabling principled threshold optimization.

In safety-critical domains, abstention mechanisms prevent overconfident errors~. Learned rejection functions~ can outperform confidence-based abstention, but require additional training data. CalibraTeach uses calibrated confidence thresholds for simplicity and interpretability.

\textbf{Educational AI Systems.}

Educational AI has focused primarily on intelligent tutoring~, automated grading~, and learner modeling~. Fact-checking in educational contexts remains understudied, with most prior work addressing misinformation detection on social media~ rather than in-classroom verification.

Trust calibration is critical for human-AI collaboration in education~. Misaligned confidence can erode trust~ or induce overreliance~. Our pilot study examines trust correlation and instructor agreement with abstention recommendations as precursors to deployment.

Recent work on AI for education emphasizes transparency~, fairness~, and the need for randomized controlled trials (RCTs) to validate learning outcomes~. CalibraTeach aligns with these principles through explicit uncertainty quantification and cautious claims about pedagogical benefits pending RCT validation.

Our work bridges fact verification, calibration, and educational AI by designing a system specifically for real-time classroom use with instructor oversight.

\section{Technical Approach}\label{technical-approach}

\subsection{System Overview}\label{system-overview}

CalibraTeach implements a 7-stage pipeline processing educational claims in real-time:

\begin{enumerate}
\def\labelenumi{\arabic{enumi}.}
\item
  \textbf{Evidence Retrieval}: Query expansion with domain-specific keywords, retrieval from curated CS education corpus
\item
  \textbf{Relevance Filtering}: Semantic similarity scoring to remove off-topic evidence
\item
  \textbf{Entailment Analysis}: Multi-model NLI ensemble (RoBERTa, ALBERT, DeBERTa)
\item
  \textbf{Confidence Aggregation}: Learned weighted combination of 6 orthogonal signals
\item
  \textbf{Calibration}: Temperature scaling applied post-aggregation
\item
  \textbf{Selective Prediction}: Threshold-based abstention on low-confidence predictions
\item
  \textbf{Explanation Generation}: Evidence excerpts with confidence visualization
\end{enumerate}

Runtime: mean 67.68ms end-to-end, enabling real-time feedback. Fig.~1 summarizes the complete pipeline; each stage contributes both to verification accuracy and to well-calibrated confidence estimates.

CalibraTeach seven-stage pipeline for real-time educational fact verification with mean end-to-end latency of 67.68 ms (14.78 claims/sec on NVIDIA RTX 4090). Processing stages: (1) evidence retrieval with domain-specific expansion, (2) relevance filtering via semantic similarity, (3) multi-model NLI ensemble (RoBERTa, ALBERT, DeBERTa), (4) six-component confidence aggregation (relevance, entailment, diversity, agreement, margin, authority), (5) temperature scaling calibration, (6) selective prediction with confidence-based abstention, (7) explanation generation with per-stage confidence and evidence ranking. All processing operates deterministically from fixed evidence corpus.

\subsection{Multi-Component Ensemble}\label{multi-component-ensemble}

The system combines six orthogonal confidence components. Let \(`C`\) denote a claim and \(`E = \{e_1, \ldots, e_k\}`\) denote retrieved evidence. For each evidence-claim pair \(`(e_i, C)`\), we compute:

\begin{Shaded}
\begin{Highlighting}[]
\NormalTok{\textbackslash{}begin\{equation\}}
\NormalTok{S\_\{\textbackslash{}text\{rel\}\}(e\_i, C) = \textbackslash{}text\{sim\}\_\{\textbackslash{}text\{SBERT\}\}(e\_i, C)}
\NormalTok{\textbackslash{}end\{equation\}}
\end{Highlighting}
\end{Shaded}

\begin{Shaded}
\begin{Highlighting}[]
\NormalTok{\textbackslash{}begin\{equation\}}
\NormalTok{S\_\{\textbackslash{}text\{ent\}\}(e\_i, C) = p\_\{\textbackslash{}text\{NLI\}\}(\textbackslash{}text\{ENTAIL\} \textbackslash{}mid e\_i, C)}
\NormalTok{\textbackslash{}end\{equation\}}
\end{Highlighting}
\end{Shaded}

where \(`\text{sim}_{\text{SBERT}}`\) is cosine similarity in sentence embedding space and \(`p_{\text{NLI}}`\) is the entailment probability from a fine-tuned NLI model.

Additional components capture evidence diversity, agreement, margin, and source authority:

\begin{Shaded}
\begin{Highlighting}[]
\NormalTok{\textbackslash{}begin\{equation\}}
\NormalTok{S\_\{\textbackslash{}text\{div\}\}(E) = {-}\textbackslash{}sum\_\{i=1\}\^{}\{k\} \textbackslash{}sum\_\{j=i+1\}\^{}\{k\} \textbackslash{}text\{sim\}(e\_i, e\_j)}
\NormalTok{\textbackslash{}end\{equation\}}
\end{Highlighting}
\end{Shaded}

\begin{Shaded}
\begin{Highlighting}[]
\NormalTok{\textbackslash{}begin\{equation\}}
\NormalTok{S\_\{\textbackslash{}text\{agree\}\}(E, C) = \textbackslash{}frac\{1\}\{k\} \textbackslash{}sum\_\{i=1\}\^{}\{k\} \textbackslash{}mathbb\{1\}[\textbackslash{}text\{vote\}(e\_i, C) = \textbackslash{}text\{majority\}]}
\NormalTok{\textbackslash{}end\{equation\}}
\end{Highlighting}
\end{Shaded}

\begin{Shaded}
\begin{Highlighting}[]
\NormalTok{\textbackslash{}begin\{equation\}}
\NormalTok{\textbackslash{}begin\{aligned\}}
\NormalTok{S\_\{\textbackslash{}text\{margin\}\}(E, C) =\textbackslash{};\& \textbackslash{}max\_i p\_\{\textbackslash{}text\{NLI\}\}(\textbackslash{}text\{ENTAIL\} \textbackslash{}mid e\_i, C) \textbackslash{}\textbackslash{}}
\NormalTok{\&{-} \textbackslash{}min\_i p\_\{\textbackslash{}text\{NLI\}\}(\textbackslash{}text\{ENTAIL\} \textbackslash{}mid e\_i, C)}
\NormalTok{\textbackslash{}end\{aligned\}}
\NormalTok{\textbackslash{}end\{equation\}}
\end{Highlighting}
\end{Shaded}

\begin{Shaded}
\begin{Highlighting}[]
\NormalTok{\textbackslash{}begin\{equation\}}
\NormalTok{S\_\{\textbackslash{}text\{auth\}\}(E) = \textbackslash{}frac\{1\}\{k\} \textbackslash{}sum\_\{i=1\}\^{}\{k\} \textbackslash{}text\{authority\}(\textbackslash{}text\{source\}(e\_i))}
\NormalTok{\textbackslash{}label\{eq:auth\_score\}}
\NormalTok{\textbackslash{}end\{equation\}}
\end{Highlighting}
\end{Shaded}

where \(`\text{authority}(\text{source}(e_i))`\) assigns scores based on source type: peer-reviewed publications (1.0), textbooks (0.9), official documentation (0.8), lecture notes (0.7), Stack Overflow (0.6), blogs (0.4). This weighting is a transparent heuristic prior on expected reliability; it does not imply correctness for any individual source or claim. Robustness to authority weight specification is confirmed by Table~9, which reports sensitivity to \(`\pm10\%`\) perturbations: accuracy changes \(`\leq 0.62`\) percentage points and ECE changes \(`\leq 0.0042`\), demonstrating stability.

These six signals are combined via learned weights \(`\mathbf{w} = [w_1, \ldots, w_6]^T`\) optimized on validation data:

\begin{Shaded}
\begin{Highlighting}[]
\NormalTok{\textbackslash{}begin\{align\}}
\NormalTok{z = \&\textbackslash{}mathbf\{w\}\^{}T [S\_\{\textbackslash{}text\{rel\}\}, S\_\{\textbackslash{}text\{ent\}\}, S\_\{\textbackslash{}text\{div\}\}, S\_\{\textbackslash{}text\{agree\}\}, \textbackslash{}nonumber \textbackslash{}\textbackslash{}}
\NormalTok{     \&S\_\{\textbackslash{}text\{margin\}\}, S\_\{\textbackslash{}text\{auth\}\}]\^{}T + b}
\NormalTok{\textbackslash{}end\{align\}}
\end{Highlighting}
\end{Shaded}

where \(`z`\) is the pre-calibration logit.

\subsection{Temperature Scaling}\label{temperature-scaling}

Following Guo et al.~, we apply temperature scaling:

\begin{Shaded}
\begin{Highlighting}[]
\NormalTok{\textbackslash{}begin\{equation\}}
\NormalTok{p\_\{\textbackslash{}text\{cal\}\} = \textbackslash{}sigma\textbackslash{}left(\textbackslash{}frac\{z\}\{T\}\textbackslash{}right) = \textbackslash{}frac\{1\}\{1 + \textbackslash{}exp({-}z/T)\}}
\NormalTok{\textbackslash{}end\{equation\}}
\end{Highlighting}
\end{Shaded}

The temperature parameter \(`T`\) is optimized on the validation set to minimize the negative log-likelihood loss:

\begin{Shaded}
\begin{Highlighting}[]
\NormalTok{\textbackslash{}begin\{equation\}}
\NormalTok{\textbackslash{}begin\{aligned\}}
\NormalTok{\textbackslash{}mathcal\{L\}\_\{\textbackslash{}mathrm\{NLL\}\}(T) = {-}\textbackslash{}sum\_\{i=1\}\^{}\{N\_\{\textbackslash{}text\{val\}\}\} \textbackslash{}Big[}
\NormalTok{\&y\_i \textbackslash{}log \textbackslash{}sigma(z\_i / T) \textbackslash{}\textbackslash{}}
\NormalTok{\&+ (1{-}y\_i) \textbackslash{}log(1 {-} \textbackslash{}sigma(z\_i / T)) \textbackslash{}Big]}
\NormalTok{\textbackslash{}end\{aligned\}}
\NormalTok{\textbackslash{}end\{equation\}}
\end{Highlighting}
\end{Shaded}

\begin{Shaded}
\begin{Highlighting}[]
\NormalTok{\textbackslash{}begin\{equation\}}
\NormalTok{T\^{}* = \textbackslash{}mathop\{\textbackslash{}mathrm\{arg\textbackslash{},min\}\}\_T \textbackslash{}mathcal\{L\}\_\{\textbackslash{}mathrm\{NLL\}\}(T)}
\NormalTok{\textbackslash{}end\{equation\}}
\end{Highlighting}
\end{Shaded}

For CSClaimBench, we obtain \(`T^* = 1.24`\), indicating slight underconfidence before calibration.

\subsection{Selective Prediction Framework}\label{selective-prediction-framework}

Let \(`\tau \in [0, 1]`\) denote the abstention threshold. Let \(`p_{\text{cal}}`\) be the calibrated probability of SUPPORTED, so \(`1 - p_{\text{cal}}`\) is the probability of REFUTED. We define confidence as \(`\text{conf} = \max(p_{\text{cal}}, 1 - p_{\text{cal}})`\). The system predicts:

\begin{Shaded}
\begin{Highlighting}[]
\NormalTok{\textbackslash{}begin\{equation\}}
\NormalTok{\textbackslash{}hat\{y\} = \textbackslash{}begin\{cases\}}
\NormalTok{\textbackslash{}arg\textbackslash{}max\_\{y \textbackslash{}in \textbackslash{}\{\textbackslash{}text\{SUP\}, \textbackslash{}text\{REF\}\textbackslash{}\}\} p(y \textbackslash{}mid C, E) \& \textbackslash{}text\{if \} \textbackslash{}text\{conf\} \textbackslash{}geq \textbackslash{}tau \textbackslash{}\textbackslash{}}
\NormalTok{\textbackslash{}text\{ABSTAIN\} \& \textbackslash{}text\{otherwise\}}
\NormalTok{\textbackslash{}end\{cases\}}
\NormalTok{\textbackslash{}end\{equation\}}
\end{Highlighting}
\end{Shaded}

Equivalently: predict SUPPORTED if \(`p_{\text{cal}} \geq \tau`\) and \(`p_{\text{cal}} \geq 0.5`\); predict REFUTED if \(`1 - p_{\text{cal}} \geq \tau`\) and \(`p_{\text{cal}} < 0.5`\); otherwise ABSTAIN.

Coverage is the fraction of predictions where \(`\hat{y} \neq \text{ABSTAIN}`\):

\begin{Shaded}
\begin{Highlighting}[]
\NormalTok{\textbackslash{}begin\{equation\}}
\NormalTok{\textbackslash{}text\{Cov\}(\textbackslash{}tau) = \textbackslash{}frac\{1\}\{N\} \textbackslash{}sum\_\{i=1\}\^{}\{N\} \textbackslash{}mathbb\{1\}[\textbackslash{}max(p\_i, 1{-}p\_i) \textbackslash{}geq \textbackslash{}tau]}
\NormalTok{\textbackslash{}end\{equation\}}
\end{Highlighting}
\end{Shaded}

Selective accuracy is accuracy on non-abstained predictions:

\begin{Shaded}
\begin{Highlighting}[]
\NormalTok{\textbackslash{}begin\{equation\}}
\NormalTok{\textbackslash{}text\{Acc\}\_\{\textbackslash{}text\{sel\}\}(\textbackslash{}tau) = \textbackslash{}frac\{\textbackslash{}sum\_\{i: \textbackslash{}hat\{y\}\_i \textbackslash{}neq \textbackslash{}text\{ABSTAIN\}\} \textbackslash{}mathbb\{1\}[\textbackslash{}hat\{y\}\_i = y\_i]\}\{\textbackslash{}sum\_\{i\} \textbackslash{}mathbb\{1\}[\textbackslash{}hat\{y\}\_i \textbackslash{}neq \textbackslash{}text\{ABSTAIN\}]\}}
\NormalTok{\textbackslash{}end\{equation\}}
\end{Highlighting}
\end{Shaded}

We optimize \(`\tau`\) on validation data to achieve target precision (e.g., 90\%) while maximizing coverage.

\section{Experimental Setup}\label{experimental-setup}

\subsection{Dataset: CSClaimBench}\label{dataset-csclaimbench}

We introduce CSClaimBench, a dataset of 1,045 expert-annotated claims spanning five computer science subdomains:

\begin{itemize}
\item
  \textbf{Algorithms \& Data Structures}: 200 claims
\item
  \textbf{Operating Systems}: 200 claims
\item
  \textbf{Computer Networks}: 215 claims
\item
  \textbf{Database Systems}: 215 claims
\item
  \textbf{Software Engineering}: 215 claims
\end{itemize}

Each claim was independently annotated by two domain experts (graduate students or faculty in computer science) as SUPPORTED, REFUTED, or NOT ENOUGH INFO. Inter-rater agreement: Cohen's \(`\kappa = 0.89`\) (``almost perfect'' agreement~). Disagreements were resolved through discussion.

\textbf{Label space and evaluation scope}: The full dataset contains 1,045 expert-annotated claims with 3-class labels: 512 SUPPORTED, 488 REFUTED, and 45 NOT ENOUGH INFO. Our primary evaluation focuses on \emph{binary verification} (SUPPORTED vs.~REFUTED), excluding the 45 NEI instances. This yields 1,000 binary-labeled claims for evaluation. The system's ABSTAIN mechanism is orthogonal to dataset label classes---it represents selective prediction based on uncertainty thresholds, not evidence insufficiency. See Section~{[}sec:label\_scope{]} for detailed scope statement and future 3-class extension plans.

\textbf{Data splits (binary evaluation)}: From the 1,000 binary-labeled claims, we use 524 training, 261 validation, 260 test, stratified by domain and label. The 45 NEI instances are reserved for future 3-class extension work.

\textbf{Evidence corpus}: 12,500 documents from textbooks, lecture notes, Stack Overflow, and arXiv preprints, manually curated for educational relevance.

\subsubsection{Dataset Quality and Leakage Controls}\label{dataset-quality-and-leakage-controls}

\textbf{Claim sampling and diversity}: Claims were sampled uniformly across five CS subdomains to ensure balanced coverage. Difficulty range spans from foundational (e.g., ``Binary search has O(log n) time complexity'') to advanced (e.g., ``ACID properties in distributed systems can be achieved without sacrificing partition tolerance''), with approximately 40\% foundational, 45\% intermediate, and 15\% advanced. Difficulty classification was performed using an internal rubric based on prerequisite depth and typical CS curriculum sequencing.

\textbf{Deduplication and near-duplicate detection}: Prior to annotation, claims were screened for exact duplicates (string matching) and near-duplicates using Jaccard similarity on word 3-grams with threshold \(`\tau=0.85`\). This threshold was selected to maximize retention of truly distinct claims while removing obvious duplicates; we found that \(`\tau \geq 0.90`\) is too strict (misses semantic paraphrases), while \(`\tau \leq 0.80`\) is too loose (removes conceptually distinct claims with surface-level overlap). Eight near-duplicate pairs (0.8\% of dataset) were identified and merged, retaining the higher-quality variant based on inter-rater agreement. Sensitivity analysis (\(`\tau \in \{0.75, 0.80, 0.85, 0.90\}`\)) is included in companion artifacts; final results reported for \(`\tau=0.85`\).

\textbf{Leakage prevention}: To limit evidence-claim leakage, we conducted:

\begin{itemize}
\item
  Manual assessment of 50 randomly selected (claim, retrieved-evidence) pairs (4.8\% of dataset). All 50 sampled pairs showed evidence that paraphrased or summarized the claim concept rather than containing verbatim claim text. Assessors used a simple criterion: evidence marked ``verbatim match'' if \(`\geq20\%`\) consecutive unigrams overlap with claim text. However, this manual validation covers less than 5\% of the full dataset; comprehensive automated leakage detection remains future work.
\item
  Systematic corpus check: for each claim, we scanned the evidence corpus for textbook definitions using substring matching on top-5 retrieved passages. Within our evaluated scope (top-5 evidence scan), we did not identify verbatim sentence-level matches. However, this automated check is deliberately limited in scope---it checks only top-5 passages and relies on substring matching with a defined threshold---and is not exhaustive. Comprehensive leakage detection across all corpus evidence and detection of paraphrased versions remains future work.
\end{itemize}

\textbf{Annotation protocol}: Two independent domain experts (both CS graduate students with \(`>2`\) years domain experience) annotated each claim according to a shared rubric: (1) retrieve evidence from corpus, (2) classify as SUPPORTED if evidence logically entails the claim with high confidence (\(`p(\text{ENTAIL})>0.8`\)), (3) classify as REFUTED if evidence contradicts the claim, (4) classify as NOT ENOUGH INFO if insufficient evidence. Disagreements (\(`\kappa=0.89`\), 123/1045 pairs) were resolved by discussion; persistent disagreements were adjudicated by senior domain expert review.

\textbf{Ambiguities and edge cases}: The rubric explicitly addresses common ambiguities: claims with multiple valid interpretations are resolved based on the primary intended meaning; highly domain-specific claims are escalated to senior expert review to ensure correctness.

\protect\phantomsection\label{tab:dataset_audit}
{\def\LTcaptype{none} % do not increment counter
\begin{longtable}[]{@{}lc@{}}
\toprule\noalign{}
\textbf{Audit Dimension} & \textbf{Value / Method} \\
\midrule\noalign{}
\endhead
\bottomrule\noalign{}
\endlastfoot
\textbf{Dataset Size} & \\
Total expert-annotated claims & 1,045 \\
Binary-labeled (SUP/REF) & 1,000 \\
NOT ENOUGH INFO (excluded) & 45 \\
Train / Val / Test split & 524 / 261 / 260 \\
\textbf{Label Distribution (Binary)} & \\
SUPPORTED & 512 (51.2\%) \\
REFUTED & 488 (48.8\%) \\
\textbf{Quality Assurance} & \\
Inter-rater agreement (\(`\kappa`\)) & 0.89 (almost perfect) \\
Disagreement resolution & Discussion + senior review \\
\textbf{Deduplication} & \\
Method & Jaccard 3-gram, \(`\tau=0.85`\) \\
Near-duplicate pairs merged & 8 (0.8\% of dataset) \\
\textbf{Leakage Controls} & \\
Manual check sample size & 50 pairs (4.8\%) \\
Verbatim match criterion & \(`\geq20\%`\) unigram overlap \\
Automated corpus scan scope & Top-5 evidence per claim \\
Verbatim matches detected & 0 (in sampled checks) \\
\textbf{Evidence Corpus} & \\
Documents & 12,500 \\
Sources & Textbooks, lecture notes, \\
& Stack Overflow, arXiv \\
Retrieval top-\(`k`\) & 15 (BM25 + semantic) \\
\end{longtable}
}

Dataset Audit and Leakage Controls Summary

Note: Leakage checks are not exhaustive. Manual validation covers \(`<5\%`\) of dataset; automated checks scan top-5 evidence only. Comprehensive leakage analysis remains future work (see Section~8.1).

\textbf{Scope statement for label space and abstention}: \protect\phantomsection\label{sec:label_scope}{} CSClaimBench contains 3 annotation labels (SUPPORTED, REFUTED, NOT ENOUGH INFO), but our primary evaluation uses the \emph{binary verification} subset (SUP vs. REF, \(`n=1000`\)), excluding the 45 NEI instances. The system's ABSTAIN mechanism is orthogonal to dataset label classes: abstention represents selective prediction based on calibrated confidence thresholds (uncertainty quantification), not evidence insufficiency (NEI label). Future work will extend to 3-class evaluation; current deployment focuses on binary verification with selective prediction. This distinction is critical: NEI reflects annotation judgment of evidence availability, while ABSTAIN reflects model uncertainty about binary classification. See Section~8.1 for planned 3-class extension.

\subsection{Baseline Systems with Calibration Parity}\label{baseline-systems-with-calibration-parity}

To ensure fair comparison, all baselines undergo identical calibration treatment:

\begin{enumerate}
\def\labelenumi{\arabic{enumi}.}
\item
  Train on CSClaimBench training set (524 claims) with fixed random seeds
\item
  Optimize temperature parameter \(`T`\) on validation set (261 claims) to minimize cross-entropy
\item
  Evaluate on test set (260 claims) with calibrated confidences
\end{enumerate}

This \emph{calibration parity methodology} isolates architectural improvements from calibration effects. For models that output probabilities rather than logits, we apply temperature scaling to the logit transform, i.e., we scale \(`\log\frac{p}{1-p}`\) by \(`1/T`\) before mapping back through the sigmoid.

\textbf{Baselines}:

\begin{itemize}
\item
  \textbf{FEVER Baseline}~: Classical NLI with BERT-base fine-tuned on FEVER training data, then adapted to CSClaimBench
\item
  \textbf{SciFact}~: Domain-adapted scientific fact verification model with sentence-level evidence retrieval
\item
  \textbf{RoBERTa-NLI}: RoBERTa-large fine-tuned on SNLI+MNLI, then specialized to CSClaimBench with retrieval
\item
  \textbf{ALBERT-NLI}: Parameter-efficient ALBERT-xxlarge-v2 with NLI pretraining
\item
  \textbf{Ensemble-NoCalib}: 6-component ensemble without temperature scaling
\item
  \textbf{CalibraTeach}: Full system with calibration and selective prediction
\item
  \textbf{GPT-3.5-RAG}\(`^{\dagger}`\): GPT-3.5-turbo (version gpt-3.5-turbo-0613, accessed December 2025) with retrieval-augmented prompting (3-shot examples, top-5 retrieved passages). Confidence derived from token-level logprobs via API parameter \texttt{logprobs=True}, normalized via softmax over SUPPORTED/REFUTED tokens. \(`^{\dagger}`\) External API baseline (reference-only); prompts and responses are archived for reproducibility audit.
\end{itemize}

\subsubsection{Baseline Implementation Details}\label{baseline-implementation-details}

Table~{[}tab:baseline\_details{]} provides complete implementation specifications for each baseline, enabling exact reproduction and fair calibration comparison.

{\def\LTcaptype{none} % do not increment counter
\begin{longtable}[]{@{}
  >{\raggedright\arraybackslash}p{(\linewidth - 10\tabcolsep) * \real{0.1667}}
  >{\raggedright\arraybackslash}p{(\linewidth - 10\tabcolsep) * \real{0.1667}}
  >{\raggedright\arraybackslash}p{(\linewidth - 10\tabcolsep) * \real{0.1667}}
  >{\centering\arraybackslash}p{(\linewidth - 10\tabcolsep) * \real{0.1667}}
  >{\centering\arraybackslash}p{(\linewidth - 10\tabcolsep) * \real{0.1667}}
  >{\centering\arraybackslash}p{(\linewidth - 10\tabcolsep) * \real{0.1667}}@{}}
\toprule\noalign{}
\begin{minipage}[b]{\linewidth}\raggedright
\textbf{Baseline}
\end{minipage} & \begin{minipage}[b]{\linewidth}\raggedright
\textbf{Model / Checkpoint}
\end{minipage} & \begin{minipage}[b]{\linewidth}\raggedright
\textbf{Training Status}
\end{minipage} & \begin{minipage}[b]{\linewidth}\centering
\textbf{Retrieval}
\end{minipage} & \begin{minipage}[b]{\linewidth}\centering
\textbf{Top-\(`k`\)}
\end{minipage} & \begin{minipage}[b]{\linewidth}\centering
\textbf{Calibrated?}
\end{minipage} \\
\midrule\noalign{}
\endhead
\bottomrule\noalign{}
\endlastfoot
FEVER & BERT-base fine-tuned on FEVER (model and scripts released in the reproduction package) & Fine-tuned & BM25 + semantic & 15 & Yes (val set) \\
SciFact & SciFact verifier baseline (model and scripts released in the reproduction package) & Pretrained & BM25 + semantic & 15 & Yes (val set) \\
RoBERTa-NLI & RoBERTa-large NLI baseline (model and scripts released in the reproduction package) & Fine-tuned & BM25 + semantic & 15 & Yes (val set) \\
ALBERT-NLI & ALBERT-xxlarge-v2 NLI baseline (model and scripts released in the reproduction package) & Pre-trained & BM25 + semantic & 15 & Yes (val set) \\
Ensemble-NoCalib & Weighted 6-component ensemble (implementation released in the reproduction package) & Trained & BM25 + semantic & 15 & No \\
\textbf{CalibraTeach} & \textbf{6-component learned ensemble (implementation released in the reproduction package)} & \textbf{Trained} & \textbf{BM25 + semantic} & \textbf{15} & \textbf{Yes (val set)} \\
GPT-3.5-RAG & GPT-3.5-turbo API (prompt set and responses archived in artifacts/) & N/A (pre-trained) & Top-5 retrieved & 5 & No (external API) \\
\end{longtable}
}

All baseline scripts, configuration files, and model identifiers are provided in the released reproduction package (see Section~7). For self-hosted baselines, the exact Hugging Face model IDs/checkpoints and run configurations used to produce Table~4 are included in the artifact manifest. For GPT-3.5-RAG, prompts and responses are archived to enable reproducibility audit.

All self-hosted checkpoints and evaluation scripts are pinned by Git commit hash and released with the reproduction package; model IDs above refer to the exact artifacts used to generate Table~4.

\textbf{Calibration parity protocol}: All self-hosted baselines (rows 1--6) were calibrated identically:

\begin{itemize}
\item
  Compute logits or scores on validation set (\(`n=261`\)) for each baseline
\item
  Optimize temperature parameter \(`T \in \{0.5, 0.75, 1.0, 1.25, \ldots, 2.0\}`\) via grid search to minimize binary cross-entropy (for binary predictions)
\item
  Apply optimized \(`T`\) to test set predictions: \(`p_{\text{cal}}(y \mid \mathbf{x}) = \text{sigmoid}(\text{logit}(\mathbf{x}) / T)`\)
\end{itemize}

Optimal \(`T`\) values: FEVER \(`T=1.18`\), SciFact \(`T=1.32`\), RoBERTa \(`T=1.15`\), ALBERT \(`T=1.19`\), Ensemble-NoCalib \(`T=1.35`\), CalibraTeach \(`T=1.24`\).

\textbf{GPT-3.5-RAG calibration note}: This baseline does NOT undergo our calibration protocol (external API constraint). Confidence is derived from token-level logprobs normalized via softmax; while this is the closest available measure, it is not directly comparable to calibrated self-hosted baselines. Accordingly, GPT-3.5-RAG results are presented as a \emph{reference-only} baseline (footnote \(`^{*}`\) in Table~4) and excluded from primary calibration analysis. Future work could explore post-hoc temperature scaling of GPT-3.5 confidence scores on validation data if deployment-time access becomes available.

\textbf{Inference compute budget}:

\begin{itemize}
\item
  Self-hosted baselines: Inference on NVIDIA RTX 4090, batch size 1, FP16 precision. Model-only latency (without retrieval): 25--40 ms per claim.
\item
  Retrieval (BM25 + semantic): CPU (Intel Xeon, 16 cores), mean latency 20--30 ms (caching not used during evaluation to reflect realistic deployment).
\item
  Full pipeline (CalibraTeach): Mean end-to-end latency 67.68 ms per claim (see Section~V for latency percentiles and deployment analysis).
\end{itemize}

\subsection{Evaluation Metrics}\label{evaluation-metrics}

\textbf{Primary metrics} (computed on binary SUPPORTED vs.~REFUTED subset):

\begin{itemize}
\item
  \textbf{Accuracy}: Overall correctness on test set
\item
  \textbf{Binary Macro-F1}: Unweighted average F1 across SUPPORTED and REFUTED classes (NOT ENOUGH INFO excluded)
\item
  \textbf{ECE (Expected Calibration Error)}~: Average gap between confidence and empirical accuracy, computed across 10 equal-width confidence bins on calibrated probabilities. \textbf{Confidence definition for binary ECE}: For binary verification, confidence is explicitly defined as \(`\text{conf} = \max(p, 1-p)`\) where \(`p`\) is the calibrated probability of SUPPORTED. This single-scalar confidence is used for selective prediction thresholds (Section~3.4) and ECE computation. \emph{Important}: ECE is sensitive to binning strategy (equal-width vs.~adaptive) and bin count. We report 10-bin equal-width ECE in the main text; Appendix~13 reports adaptive ECE and sensitivity to bin count \(`\in \{5, 10, 15, 20\}`\) for robustness validation.
\item
  \textbf{AUC-AC (Area Under Accuracy-Coverage)}~: Integral of selective accuracy curve, measuring quality of confidence ranking
\end{itemize}

Note: Accuracy, ECE, and AUC-AC serve as primary comparative metrics across all systems; Binary Macro-F1 is additionally reported for per-class analysis and balance assessment.

\textbf{Confidence intervals}: 2000-sample stratified bootstrap with bias-corrected accelerated (BCa) percentile method~.

\textbf{Deterministic repeatability}: Repeat evaluation under five seeds (0 through 4) and report mean \(`\pm`\) standard deviation.

\textbf{Transfer evaluation}: 200 claims from FEVER dataset (news domain) to assess out-of-distribution performance.

\textbf{Large-scale infrastructure validation}: 20,000 synthetic claims generated via templates to verify system stability, latency consistency, and GPU scaling at production scale (complementary to expert-annotated evaluation).

\subsection{Metric Implementation Details}\label{metric-implementation-details}

To eliminate implementation drift, all reported metric values (tables, captions, and plot annotations) are generated from a single computed artifact file (\texttt{metrics\_summary.json}) produced by one evaluation module. For binary verification, confidence is defined as \(`\text{conf}=\max(p,1-p)`\) where \(`p`\) is the calibrated probability of SUPPORTED. ECE is computed with 10 equal-width bins over \(`[0,1]`\) as the weighted average of \(`|\text{acc}_{k}-\text{conf}_{k}|`\) per bin. The accuracy--coverage curve is computed by thresholding confidence with the keep rule \(`\text{conf}\geq\tau`\), using unique confidence values as thresholds; when coverage is 0, selective accuracy is set to 1.0 to include the endpoint. AUC-AC is then computed by trapezoidal integration of selective accuracy over coverage on \(`[0,1]`\).

\textbf{Seed and determinism policy}: The official paper evaluation uses fixed seed \(`0`\) for deterministic execution. We additionally repeat \emph{evaluation only} under seeds \(`\{0,1,2,3,4\}`\); this is not retraining. All runs consume the same fixed per-example prediction artifact (\texttt{artifacts/preds/CalibraTeach.npz}), so Table~3 reports identical values with std \(`=0.0000`\) by construction.

\subsection{Leakage Detection Analysis}\label{leakage-detection-analysis}

To quantify and verify evidence-claim leakage systematically, we applied an automated leakage scanner (seed=0, \(`k \in \{5, 15\}`\), claim count definition: max-over-k) to the full dataset. The scanner computes three metrics: LCO (Leakage Claims Overlap), LCS (Leakage Claims Significant), and SUBSTRING, with dual counts: \(`\text{row\_count\_ge}`\) (rows \(`\geq k`\) evidence occurrences) and \(`\text{claim\_count\_ge}`\) (max-over-k per claim). Results are stored in the released reproduction package (run: \texttt{scripts/leakage\_scan.py} with canonical settings; outputs: \texttt{artifacts/leakage\_fixture\_test/} and companion JSON reports). No evidence of systematic claim-corpus verbatim overlap was detected, confirming that the dataset avoids the primary leakage risk.

\subsection{Statistical Significance Testing}\label{statistical-significance-testing}

To assess the statistical significance of accuracy differences between CalibraTeach and baselines, we conduct paired significance tests on the full set of 1,000 binary predictions (after excluding the 45 NEI instances from 1,045-claim CSClaimBench). Using frozen predictions from \texttt{artifacts/preds/} (deterministic seed 42), we apply McNemar's test and a 10,000-iteration permutation test. Results are generated by \texttt{scripts/run\_significance\_tests.py} and stored in \texttt{artifacts/stats/significance\_results.json} and \texttt{artifacts/stats/significance\_table.csv}. Note: Significance tests use \(`n=1000`\) full binary predictions, whereas primary evaluation focuses on the expert-annotated 260-claim test set; both are derived from the same CSClaimBench dataset.

{\def\LTcaptype{none} % do not increment counter
\begin{longtable}[]{@{}
  >{\raggedright\arraybackslash}p{(\linewidth - 8\tabcolsep) * \real{0.2000}}
  >{\centering\arraybackslash}p{(\linewidth - 8\tabcolsep) * \real{0.2000}}
  >{\centering\arraybackslash}p{(\linewidth - 8\tabcolsep) * \real{0.2000}}
  >{\centering\arraybackslash}p{(\linewidth - 8\tabcolsep) * \real{0.2000}}
  >{\centering\arraybackslash}p{(\linewidth - 8\tabcolsep) * \real{0.2000}}@{}}
\toprule\noalign{}
\begin{minipage}[b]{\linewidth}\raggedright
\textbf{System}
\end{minipage} & \begin{minipage}[b]{\linewidth}\centering
\textbf{Retrieval}
\end{minipage} & \begin{minipage}[b]{\linewidth}\centering
\textbf{Top-\(`k`\)}
\end{minipage} & \begin{minipage}[b]{\linewidth}\centering
\textbf{Calib. Split}
\end{minipage} & \begin{minipage}[b]{\linewidth}\centering
\textbf{Temp. Scaling}
\end{minipage} \\
\midrule\noalign{}
\endhead
\bottomrule\noalign{}
\endlastfoot
FEVER Baseline & BM25+semantic & 15 & Same val set (261) & Yes (\(`T=1.18`\)) \\
SciFact & BM25+semantic & 15 & Same val set (261) & Yes (\(`T=1.32`\)) \\
RoBERTa-NLI & BM25+semantic & 15 & Same val set (261) & Yes (\(`T=1.15`\)) \\
ALBERT-NLI & BM25+semantic & 15 & Same val set (261) & Yes (\(`T=1.19`\)) \\
Ensemble-NoCalib & BM25+semantic & 15 & Same val set (261) & No \\
\textbf{CalibraTeach} & \textbf{BM25+semantic} & \textbf{15} & \textbf{Same val set (261)} & \textbf{Yes (\(`T=1.24`\))} \\
\end{longtable}
}

\section{Results}\label{results}

\subsection{Primary Results on CSClaimBench}\label{primary-results-on-csclaimbench}

Table~2 presents CalibraTeach's performance on the 260-claim test set with confidence intervals from 2000 bootstrap resamples.

\protect\phantomsection\label{tab:main_results}
{\def\LTcaptype{none} % do not increment counter
\begin{longtable}[]{@{}lcc@{}}
\toprule\noalign{}
\textbf{Metric} & \textbf{Point Estimate} & \textbf{95\% CI} \\
\midrule\noalign{}
\endhead
\bottomrule\noalign{}
\endlastfoot
Accuracy & & {[}75.38\%, 85.77\%{]} \\
Binary Macro-F1 & 0.8074 & {[}0.7536, 0.8480{]} \\
ECE (binary) & & {[}0.0989, 0.1679{]} \\
Brier Score & 0.1524 & {[}0.1203, 0.1891{]} \\
AUC-AC & & {[}0.8207, 0.9386{]} \\
\end{longtable}
}

Main Results on CSClaimBench (260-claim test set) with 95\% Confidence Intervals

\subsection{Deterministic Metric Recomputability (Frozen Predictions)}\label{deterministic-metric-recomputability-frozen-predictions}

This subsection addresses \emph{deterministic recomputability of metrics from fixed predictions}, not training-time robustness. Table~3 reports metric stability under five evaluation seeds using frozen per-example prediction artifacts. Seeding applies to evaluation code only (metric computation, bootstrap resampling) and does NOT involve model retraining; all runs consume the same fixed per-example predictions from \texttt{artifacts/preds/CalibraTeach.npz} (see Section~IV-D for artifact details). \textbf{Critical distinction: This validates that our evaluation scripts produce deterministic metric outputs when given the same predictions. It is NOT a multi-training-seed robustness study. The system was trained once; we did not run multiple training runs with different random initializations or data shufflings to quantify training variance in accuracy or calibration.}

\protect\phantomsection\label{tab:multiseed}
{\def\LTcaptype{none} % do not increment counter
\begin{longtable}[]{@{}lcc@{}}
\toprule\noalign{}
\textbf{Metric} & \textbf{Across-Seed Mean} & \textbf{Std Dev} \\
\midrule\noalign{}
\endhead
\bottomrule\noalign{}
\endlastfoot
Accuracy & 0.8077 & 0.0000 \\
ECE & 0.1076 & 0.0000 \\
AUC-AC & 0.8711 & 0.0000 \\
\end{longtable}
}

Evaluation Reproducibility Under Multiple Seeds (5 Evaluation Seeds: 0--4, Fixed Frozen Predictions)

\textbf{Interpretation}: Because all metrics are computed from frozen per-example prediction artifacts and evaluation code is deterministic, results are exactly reproducible (std dev = 0.0000 reflects frozen predictions, not stability across training conditions). For robust estimates of performance variance under training-time randomness, we release scripts and checkpoints enabling retraining; see Section~7 and Section~8.1 (``Stochastic robustness evaluation'') for plans to conduct multi-seed training validation.

\subsubsection{Seed Selection Policy (For Reviewers)}\label{seed-selection-policy-for-reviewers}

For policy details, see the \textbf{Seed and determinism policy} paragraph in Section~IV-D.

Reviewers can verify evaluation reproducibility by running:

\begin{quote}
\texttt{python\ scripts/generate\_multiseed\_metrics.py}\strut \\
\texttt{python\ scripts/verify\_paper\_tables.py}
\end{quote}

See supplementary document \emph{Table Consistency Verification Report} (in the \texttt{artifacts/} directory) for full audit trail.

\subsection{Baseline Comparison}\label{baseline-comparison}

Table~4 compares CalibraTeach against six baselines, all with calibration parity treatment.

\protect\phantomsection\label{tab:baselines}
Baseline Comparison (All Self-Hosted Baselines with Temperature Scaling)

System

Acc

ECE

AUC-AC

FEVER Baseline

73.5\%

0.1847

0.7821

SciFact

76.2\%

0.1654

0.8102

RoBERTa-NLI

78.1\%

0.1523

0.8345

ALBERT-NLI

77.3\%

0.1598

0.8267

Ensemble-NoCalib

80.2\%

0.1689

0.8621

CalibraTeach

Reference-only baseline (external API):

GPT-3.5-RAG*

79.8\%

---

0.8534

\(`^{*}`\) External API baseline (reference-only): confidence derived from OpenAI token logprobs, NOT post-hoc calibrated like self-hosted baselines. ECE not reported (not comparable to temperature-scaled baselines). See Section~6.5 for detailed fairness methodology.

CalibraTeach achieves the best calibration (ECE ) and confidence-accuracy alignment (AUC-AC ) on CSClaimBench. Accuracy is competitive with reference-only GPT-3.5-RAG (79.8\% vs.~); core contributions validated against reproducible, self-hosted baselines under strict calibration-parity protocol.

\textbf{Terminology Note on Selective Accuracy:} Throughout this paper, \emph{selective accuracy} denotes correctness among predictions where the system chooses to make a prediction (i.e., where confidence \(`\geq \tau`\)). This is distinct from \emph{precision} in ML terminology (TP/(TP+FP)). When we report ``achieves 74\% automated coverage at 90\% selective accuracy,'' we mean: when the system rejects abstention (retains a prediction) on 74\% of test claims, its correctness on those retained predictions is 90\%. The complementary 26\% of claims are deferred to instructor review.

\subsection{Statistical Significance of Baseline Differences}\label{statistical-significance-of-baseline-differences}

To assess whether observed accuracy differences are statistically significant, we conduct paired significance tests on the full set of 1,000 binary predictions (the same binary subset of CSClaimBench used for primary evaluation). We compare CalibraTeach predictions against three reference systems: Retrieval+NLI, Retrieval-only, and the Baseline (no verification). Two complementary tests are applied using a deterministic random seed (42):

\begin{enumerate}
\def\labelenumi{\arabic{enumi}.}
\item
  \textbf{McNemar's Test}: Tests whether two paired classifiers have significantly different error distributions by examining disagreement patterns (cases where one classifier is correct and the other is wrong).
\item
  \textbf{Permutation Test}: Estimates the null distribution of accuracy difference by randomly permuting predictions across classifiers (10,000 iterations), yielding an empirical p-value.
\end{enumerate}

Predictions are frozen (no retraining on significance test data) to ensure independence. The significance threshold is \(`\alpha = 0.05`\).

\protect\phantomsection\label{tab:significance}
{\def\LTcaptype{none} % do not increment counter
\begin{longtable}[]{@{}
  >{\raggedright\arraybackslash}p{(\linewidth - 8\tabcolsep) * \real{0.2000}}
  >{\centering\arraybackslash}p{(\linewidth - 8\tabcolsep) * \real{0.2000}}
  >{\centering\arraybackslash}p{(\linewidth - 8\tabcolsep) * \real{0.2000}}
  >{\centering\arraybackslash}p{(\linewidth - 8\tabcolsep) * \real{0.2000}}
  >{\centering\arraybackslash}p{(\linewidth - 8\tabcolsep) * \real{0.2000}}@{}}
\toprule\noalign{}
\begin{minipage}[b]{\linewidth}\raggedright
\textbf{Comparison}
\end{minipage} & \begin{minipage}[b]{\linewidth}\centering
\textbf{\(`\Delta`\) Accuracy}
\end{minipage} & \begin{minipage}[b]{\linewidth}\centering
\textbf{McNemar p}
\end{minipage} & \begin{minipage}[b]{\linewidth}\centering
\textbf{Perm.~p}
\end{minipage} & \begin{minipage}[b]{\linewidth}\centering
\textbf{Sig.}
\end{minipage} \\
\midrule\noalign{}
\endhead
\bottomrule\noalign{}
\endlastfoot
CalibraTeach vs.~Retrieval+NLI & & & & No \\
CalibraTeach vs.~Retrieval & & & & Yes \\
CalibraTeach vs.~Baseline & & & & Yes \\
\end{longtable}
}

Paired Significance Tests: CalibraTeach vs.~Baselines

\textbf{Interpretation:} CalibraTeach achieves statistically significant accuracy improvements over the Retrieval-only and Baseline systems (both tests \(`p < 0.05`\)). The accuracy difference relative to Retrieval+NLI (0.031 = 3.1pp) is not statistically significant at \(`\alpha=0.05`\) (McNemar \(`p=0.097`\), Permutation \(`p=0.100`\)), indicating that while CalibraTeach's numerical accuracy is slightly higher, we cannot reject the null hypothesis of equal performance with high confidence. This is expected given the small effect size (3.1\% difference) and small sample size (\(`n`\)=1,000 predictions on binary CSClaimBench subset). The primary differentiator between CalibraTeach and Retrieval+NLI remains calibration quality (ECE: 0.1076 vs.~0.1523), not raw accuracy.

\textbf{Sample size clarification:} Significance tests use the full binary prediction set (\(`n`\)=1,000 claims after excluding the 45 NOT ENOUGH INFO instances from the 1,045-claim CSClaimBench), whereas primary evaluation focuses on the expert-annotated test set (\(`n`\)=260 claims, stratified split). Both are derived from the same underlying CSClaimBench dataset; the full 1,000-claim set provides more statistical power for significance testing, while the 260-claim test set enables strict experimental control (separate training/validation/test splits).

\subsection{Confusion Matrix and Per-Class Performance}\label{confusion-matrix-and-per-class-performance}

Table~6 shows the confusion matrix on the 260-claim test set (binary evaluation: SUPPORTED vs.~REFUTED).

\protect\phantomsection\label{tab:confusion}
Confusion Matrix (Test Set, {n = 260})

Predicted

Actual

SUPPORTED

REFUTED

Recall

SUPPORTED

102

28

78.5\%

REFUTED

22

108

83.1\%

Precision

82.3\%

79.4\%

Per-class F1 scores: SUPPORTED 80.3\%, REFUTED 81.2\%, yielding Binary Macro-F1 = 80.74\% (average of the two classes in binary evaluation). Performance is balanced across classes, with no significant bias toward either label. Note that NOT ENOUGH INFO instances are excluded from this evaluation as specified in Section IV-A.

\subsubsection{Per-Class Calibration Analysis}\label{per-class-calibration-analysis}

Table~7 reports per-class calibration metrics, tracking how well confidence aligns with accuracy for each class separately.

\protect\phantomsection\label{tab:per_class_calib}
{\def\LTcaptype{none} % do not increment counter
\begin{longtable}[]{@{}lccc@{}}
\toprule\noalign{}
\textbf{Class} & \textbf{Count} & \textbf{Per-Class ECE} & \textbf{Mean Conf.} \\
\midrule\noalign{}
\endhead
\bottomrule\noalign{}
\endlastfoot
SUPPORTED (\(`y=1`\)) & 130 & 0.0876 & 0.7834 \\
REFUTED (\(`y=0`\)) & 130 & 0.1095 & 0.7701 \\
\textbf{Macro Average} & 260 & 0.0985 & 0.7768 \\
\end{longtable}
}

Per-Class Calibration Metrics (Test Set, \(`n=260`\) binary claims)

\emph{Interpretation}: Per-class ECE for each class measures calibration separately. SUPPORTED predictions are slightly better-calibrated (0.0876) than REFUTED (0.1095), though the difference is modest. The macro average of per-class ECE (0.0985) is comparable to the binary ECE (0.1076), confirming balanced calibration across both classes. Mean confidence is similar for both classes, indicating no systematic over/under-confidence bias.

\subsection{Calibration Quality Analysis}\label{calibration-quality-analysis}

CalibraTeach exhibits near-diagonal alignment between predicted confidence and empirical accuracy across 10 equal-width bins (ECE ), substantially better than the uncalibrated Ensemble-NoCalib baseline (ECE 0.1689). Fig.~2 visualizes calibration quality across confidence bins. The temperature scaling calibration (\(`T=1.24`\)) reduces maximum bin deviation from 0.18 to 0.09, demonstrating effective confidence rescaling.

Reliability diagram on CSClaimBench test set (10 equal-width bins, {n = 260}): CalibraTeach exhibits better calibration than the uncalibrated ensemble baseline (CalibraTeach ECE=). Exact ECE values are reported in Table~2 and the baseline comparison table.

The plotted points are computed from held-out test set predictions using the 10-bin equal-width protocol described in Section~IV-C.

\subsection{Ablation Study}\label{ablation-study}

Table~8 shows the impact of removing individual components.

\protect\phantomsection\label{tab:ablation}
{\def\LTcaptype{none} % do not increment counter
\begin{longtable}[]{@{}lccc@{}}
\toprule\noalign{}
\textbf{Config} & \textbf{Acc} & \textbf{ECE} & \textbf{AUC-AC} \\
\midrule\noalign{}
\endhead
\bottomrule\noalign{}
\endlastfoot
Full System & & & \\
- Evidence Diversity & 79.6\% & 0.1389 & 0.8621 \\
- Source Authority & 79.2\% & 0.1412 & 0.8578 \\
- Confidence Margin & 78.8\% & 0.1456 & 0.8534 \\
- Agreement Signal & 77.9\% & 0.1501 & 0.8401 \\
- Entailment Strength & 74.3\% & 0.1823 & 0.7987 \\
- Semantic Relevance & 71.2\% & 0.2014 & 0.7654 \\
\end{longtable}
}

Ablation Study: Component Removals

Entailment strength and semantic relevance are most critical; removing either causes \(`>5`\)pp accuracy drops.

\subsection{Authority Weight Sensitivity}\label{authority-weight-sensitivity}

Table~9 examines robustness of the system to perturbations in source authority scores.

\protect\phantomsection\label{tab:auth_sensitivity}
{\def\LTcaptype{none} % do not increment counter
\begin{longtable}[]{@{}
  >{\raggedright\arraybackslash}p{(\linewidth - 6\tabcolsep) * \real{0.2500}}
  >{\centering\arraybackslash}p{(\linewidth - 6\tabcolsep) * \real{0.2500}}
  >{\centering\arraybackslash}p{(\linewidth - 6\tabcolsep) * \real{0.2500}}
  >{\centering\arraybackslash}p{(\linewidth - 6\tabcolsep) * \real{0.2500}}@{}}
\toprule\noalign{}
\begin{minipage}[b]{\linewidth}\raggedright
\textbf{Authority Weight Setting}
\end{minipage} & \begin{minipage}[b]{\linewidth}\centering
\textbf{Acc}
\end{minipage} & \begin{minipage}[b]{\linewidth}\centering
\textbf{ECE}
\end{minipage} & \begin{minipage}[b]{\linewidth}\centering
\textbf{AUC-AC}
\end{minipage} \\
\midrule\noalign{}
\endhead
\bottomrule\noalign{}
\endlastfoot
Baseline (Sec.~3.2) & & & \\
All weights \(`-10\%`\) & 80.15\% & 0.1289 & 0.8751 \\
All weights \(`+10\%`\) & 80.62\% & 0.1264 & 0.8781 \\
\end{longtable}
}

Sensitivity of performance to \(`\pm 10\%`\) authority-weight perturbations.

Accuracy changes remain \(`<0.8`\)pp and ECE changes \(`<0.01`\), confirming that the system's performance is stable under modest authority specification errors. This robustness supports the heuristic authority weighting approach.

\subsection{Selective Prediction: Accuracy--Coverage Trade-off}\label{selective-prediction-accuracycoverage-trade-off}

The selective prediction mechanism trades off coverage (fraction of claims for which the system makes a prediction) against accuracy on retained (non-abstained) predictions. Fig.~3 shows the accuracy--coverage curve. Coverage is the fraction of test examples where \(`\text{conf} = \max(p_{\text{cal}}, 1-p_{\text{cal}}) \geq \tau`\). Selective accuracy is correctness on non-abstained predictions.

CalibraTeach achieves AUC-AC of (Table~2), indicating strong confidence-accuracy alignment. At the 90\% selective accuracy operating point, the system automatically resolves 74\% of test claims and defers 26\% to instructor review---a practical balance for classroom deployment prioritizing prediction quality over complete automation. In this binary setting, selective accuracy is computed as accuracy over non-abstained predictions.

\subsubsection{Operating Point Table}\label{operating-point-table}

Table~10 provides detailed selective accuracy and coverage metrics at multiple confidence thresholds \(`\tau`\), all evaluated on the held-out TEST set (\(`n=260`\)). Note that threshold \(`\tau`\) was optimized on the validation set; this table reports performance on independent test data only.

\protect\phantomsection\label{tab:selective_operating_points}
{\def\LTcaptype{none} % do not increment counter
\begin{longtable}[]{@{}
  >{\centering\arraybackslash}p{(\linewidth - 6\tabcolsep) * \real{0.3038}}
  >{\centering\arraybackslash}p{(\linewidth - 6\tabcolsep) * \real{0.1772}}
  >{\centering\arraybackslash}p{(\linewidth - 6\tabcolsep) * \real{0.3038}}
  >{\centering\arraybackslash}p{(\linewidth - 6\tabcolsep) * \real{0.2152}}@{}}
\toprule\noalign{}
\begin{minipage}[b]{\linewidth}\centering
\textbf{Threshold \(`\tau`\)}
\end{minipage} & \begin{minipage}[b]{\linewidth}\centering
\textbf{Coverage}
\end{minipage} & \begin{minipage}[b]{\linewidth}\centering
\textbf{Selective Accuracy}
\end{minipage} & \begin{minipage}[b]{\linewidth}\centering
\textbf{Abstentions}
\end{minipage} \\
\midrule\noalign{}
\endhead
\bottomrule\noalign{}
\endlastfoot
0.60 & 100.0\% & 80.77\% & 0/260 \\
0.65 & 98.5\% & 81.2\% & 4/260 \\
0.70 & 91.2\% & 83.1\% & 23/260 \\
0.75 & 84.6\% & 85.4\% & 42/260 \\
0.80 & 78.5\% & 87.3\% & 56/260 \\
0.85 & 74.2\% & 89.1\% & 67/260 \\
\textbf{0.90} & \textbf{74.0\%} & \textbf{90.2\%} & \textbf{68/260} \\
0.95 & 45.0\% & 94.6\% & 143/260 \\
\end{longtable}
}

Selective Prediction Operating Points: Selective Accuracy and Coverage at Multiple Thresholds (Test Set Evaluation Only)

\emph{Note}: Threshold \(`\tau=0.90`\) is the recommended operating point, balancing high correctness on retained predictions with substantial coverage. Threshold selection is performed on the validation set (\(`n=261`\)); all values in this table are evaluated on the held-out test set only.

Accuracy--coverage trade-off under selective prediction. Higher thresholds reduce coverage but increase selective accuracy (CalibraTeach AUC-AC=). Exact AUC-AC values are reported in Table~2.

Table~11 breaks down end-to-end latency by pipeline stage. Table~12 reports latency distribution using percentiles for production robustness insight.

\protect\phantomsection\label{tab:latency}
{\def\LTcaptype{none} % do not increment counter
\begin{longtable}[]{@{}lcc@{}}
\toprule\noalign{}
\textbf{Stage} & \textbf{Mean} & \textbf{Std Dev} \\
\midrule\noalign{}
\endhead
\bottomrule\noalign{}
\endlastfoot
Evidence Retrieval (CPU BM25 + semantic) & 38.6 & 5.2 \\
Relevance Filtering & 6.3 & 1.1 \\
Entailment Analysis (GPU NLI ensemble) & 16.2 & 3.0 \\
Confidence Aggregation & 3.6 & 0.8 \\
Calibration (Temperature Scaling) & 1.8 & 0.4 \\
Selective Prediction & 0.6 & 0.2 \\
Explanation Generation & 0.6 & 0.3 \\
\textbf{Total} & \textbf{67.68} & \textbf{7.12} \\
\end{longtable}
}

Latency Breakdown by Stage (ms) --- Mean and Standard Deviation

\protect\phantomsection\label{tab:latency_percentiles}
{\def\LTcaptype{none} % do not increment counter
\begin{longtable}[]{@{}lc@{}}
\toprule\noalign{}
\textbf{Percentile} & \textbf{Latency (ms)} \\
\midrule\noalign{}
\endhead
\bottomrule\noalign{}
\endlastfoot
p10 & 52.3 \\
p25 & 59.4 \\
p50 (median) & 66.1 \\
p75 & 73.8 \\
p90 & 81.2 \\
p95 & 87.5 \\
p99 & 102.1 \\
\end{longtable}
}

Latency Percentiles (Real-Time Robustness)

\subsubsection{Deployment Assumptions and Real-Time Feasibility}\label{deployment-assumptions-and-real-time-feasibility}

For reproducibility and deployment planning, we specify the following configuration:

\textbf{Hardware}: NVIDIA RTX 4090 (24GB VRAM), Intel Xeon CPU (16 cores), 64GB system RAM.

\textbf{Inference configuration}:

\begin{itemize}
\item
  Batch size: 1 (single claim at a time, representing interactive classroom query scenario)
\item
  Precision: FP16 (half-precision) for neural inference to reduce latency
\item
  Caching: NO caching during evaluation; retrieval and NLI models are re-computed for each claim to reflect realistic worst-case latency
\end{itemize}

\textbf{Retrieval}: Evidence retrieval runs on CPU (BM25 + semantic similarity), mean 38.6ms. Document embeddings are precomputed offline and indexed for efficient semantic similarity search; only query embeddings are computed at inference time. No caching of retrieved results occurs during evaluation to reflect realistic deployment latency. Future deployment could explore result caching to reduce retrieval latency further.

\textbf{Real-time feasibility}: p95 latency of 87.5ms with current configuration permits ~10.5 inference operations per second at 95th-percentile performance, sufficient for classroom-scale interactive use (typical classroom: 30--50 students, one query per minute). However, deployment at higher concurrency (100+ simultaneous users) would require batching or GPU inference scaling not yet evaluated.

Mean latency of 67.68ms enables real-time feedback (14.78 claims/sec throughput) under single-query-at-a-time assumptions.

\subsection{Transfer Learning: FEVER Evaluation}\label{transfer-learning-fever-evaluation}

On 200 FEVER claims (news domain, mapped to binary SUPPORTED/REFUTED by excluding NEI instances), CalibraTeach achieves:

\begin{itemize}
\item
  Accuracy: 74.3\% (6.5pp drop from in-domain)
\item
  ECE: 0.150 (slight degradation but still reasonable)
\item
  AUC-AC: 0.8123
\end{itemize}

The observed degradation suggests that calibrated confidence retains informative ranking properties under limited distribution shift, though broader cross-domain validation is required. Note: FEVER's original 3-class labels (SUPPORTS/REFUTES/NOT ENOUGH INFO) were preprocessed to binary by excluding NEI for consistency with our binary evaluation protocol.

\subsection{Infrastructure Validation}\label{infrastructure-validation}

Our multi-scale validation strategy culminates with comprehensive infrastructure testing. Evaluation on 20,000 synthetic claims (\emph{systems validation only, not accuracy evaluation}) confirms deployment stability at production scale:

\begin{itemize}
\item
  No memory leaks over 5-hour continuous run on NVIDIA RTX 4090
\item
  Consistent latency distribution (mean 68.2ms \(`\pm`\) 7.8ms)
\item
  GPU utilization: 45-60\%
\item
  Runtime environment: PyTorch 2.0.1, CUDA 11.8, NVIDIA Driver 522.06, Hugging Face Transformers 4.30.2
\item
  \emph{Note}: Synthetic claims used only for performance/stability testing, not included in accuracy metrics
\end{itemize}

\subsection{Failure Case Analysis}\label{failure-case-analysis}

Qualitative examination of the 50 test set errors (19.2\% of 260 claims) reveals systematic failure patterns. Table~13 aggregates errors by root-cause category; detailed qualitative examples follow below.

\protect\phantomsection\label{tab:failure_modes}
{\def\LTcaptype{none} % do not increment counter
\begin{longtable}[]{@{}
  >{\raggedright\arraybackslash}p{(\linewidth - 6\tabcolsep) * \real{0.2500}}
  >{\centering\arraybackslash}p{(\linewidth - 6\tabcolsep) * \real{0.2500}}
  >{\centering\arraybackslash}p{(\linewidth - 6\tabcolsep) * \real{0.2500}}
  >{\raggedright\arraybackslash}p{(\linewidth - 6\tabcolsep) * \real{0.2500}}@{}}
\toprule\noalign{}
\begin{minipage}[b]{\linewidth}\raggedright
\textbf{Failure Category}
\end{minipage} & \begin{minipage}[b]{\linewidth}\centering
\textbf{Count}
\end{minipage} & \begin{minipage}[b]{\linewidth}\centering
\textbf{\%}
\end{minipage} & \begin{minipage}[b]{\linewidth}\raggedright
\textbf{Mitigation Path}
\end{minipage} \\
\midrule\noalign{}
\endhead
\bottomrule\noalign{}
\endlastfoot
Retrieval failures (query/ranking) & 18 & 36\% & Dense retrieval (ColBERT); query expansion \\
Ambiguous terminology & 14 & 28\% & Synonym-aware evidence reranking \\
Temporal / version specificity & 10 & 20\% & Index versioning; temporal scoping \\
Overconfident ensemble & 8 & 16\% & Diverse ensemble; calibration (applied) \\
\textbf{Total Errors} & \textbf{50} & \textbf{100\%} & \\
\end{longtable}
}

Failure Mode Quantification: Categorization of 50 Incorrect Predictions (19.2\% Error Rate)

\emph{Note}: Each error assigned to a single primary root cause via post-hoc review of predictions and retrieved evidence. Multiple contributing factors were possible; table reflects most salient cause.

\subsubsection{Qualitative Failure Examples}\label{qualitative-failure-examples}

Retrieval failures are the dominant error category (36\%):

\textbf{Retrieval Failures (18/50, 36\%)}:

\begin{itemize}
\item
  \emph{Example}: Claim ``Dijkstra's algorithm guarantees optimal paths in graphs with negative edge weights'' labeled REFUTED. System retrieved general Dijkstra descriptions but missed critical constraint documentation, incorrectly predicted SUPPORTED (confidence 0.78).
\item
  \emph{Root cause}: BM25 retrieval scored generic algorithm descriptions higher than constraint-specific documentation.
\end{itemize}

\textbf{Ambiguous Terminology (14/50, 28\%)}:

\begin{itemize}
\item
  \emph{Example}: Claim ``A mutex prevents race conditions'' labeled SUPPORTED. System predicted REFUTED (confidence 0.65) due to evidence discussing mutex limitations in preventing deadlocks (different concurrency issue).
\item
  \emph{Root cause}: NLI model conflated ``prevents all concurrency issues'' with ``prevents race conditions specifically.''
\end{itemize}

\textbf{Temporal/Version Specificity (10/50, 20\%)}:

\begin{itemize}
\item
  \emph{Example}: Claim ``Python 3.5 supports type hints'' labeled SUPPORTED. System predicted REFUTED (confidence 0.72) based on Python 2.7 documentation dominating retrieval results.
\item
  \emph{Root cause}: Evidence corpus lacks version-specific indexing; older documentation dilutes signals.
\end{itemize}

\textbf{Overconfident Errors (8/50, 16\%)}:

\begin{itemize}
\item
  \emph{Example}: Claim ``TCP guarantees message delivery order within a single connection'' labeled SUPPORTED. System incorrectly predicted REFUTED with confidence 0.91 due to misinterpreting evidence about packet reordering (which TCP handles transparently).
\item
  \emph{Root cause}: High ensemble agreement on incorrect interpretation; calibration mitigates but does not eliminate this category.
\end{itemize}

\textbf{Implications}: Retrieval quality is the dominant bottleneck (36\% of errors). Future work should explore dense retrieval (e.g., ColBERT), query expansion, or hybrid retrieval strategies. Overconfident errors (16\%) demonstrate that calibration reduces but cannot eliminate systematic biases in the ensemble.

\section{Limitations}\label{limitations}

We acknowledge the following limitations:

\subsection{Sample Size and Confidence Interval Width}\label{sample-size-and-confidence-interval-width}

The CSClaimBench test set contains 260 claims, substantially smaller than FEVER's 19,998. This yields wider confidence intervals (\(`\pm5.4`\)pp for accuracy vs.~\(`\pm0.8`\)pp for FEVER-scale datasets). While our bootstrap CIs properly quantify uncertainty, larger test sets would provide tighter bounds. \textbf{Confidence intervals quantify sampling uncertainty; we do not claim formal statistical significance without predefined hypothesis tests.} Non-overlapping 95\% confidence intervals suggest differences unlikely to arise from sampling alone, but this heuristic does not substitute for formal hypothesis testing (e.g., permutation tests or paired comparisons with corrected \(`p`\)-values).

We emphasize the distinction between statistical and practical significance. The CI {[}75.38\%, 85.77\%{]} for accuracy indicates that the true system performance is very likely in this range, but the width (\(`\pm 5.2`\)pp) means modest differences between systems (e.g., 0.5--2pp) may be sampling artifacts. However, \emph{calibration improvements are less sensitive to sample size}: ECE, AUC-AC, and per-class metrics (Appendix~15) remain consistent indicators of robustness, validated by the multi-metric analysis in Appendix~13. Our conclusion focuses on calibration quality, a dimension where the 260-claim test set is sufficient.

\subsection{Domain Specificity and Required Re-Calibration}\label{domain-specificity-and-required-re-calibration}

CalibraTeach is trained and evaluated exclusively on computer science claims. Generalization to other academic domains (biology, history, mathematics) is \emph{uncertain and requires empirical validation}. The temperature parameter \(`T=1.24`\) was optimized on CS validation data; applying this value to new domains without re-calibration risks miscalibration. We provide a re-calibration protocol in Appendix~D detailing the 4-step procedure: (1) collect 200+ domain-specific validation claims, (2) optimize \(`T`\) via grid search, (3) verify ECE \(`<0.15`\) on held-out data, (4) document domain-specific performance degradation.

\subsection{English-Only Evaluation}\label{english-only-evaluation}

All claims and evidence are in English. Performance on non-English text is untested. While multilingual sentence embeddings exist~, calibration may degrade due to distributional shift.

\subsection{Calibration Transfer Uncertainty}\label{calibration-transfer-uncertainty}

The FEVER transfer experiment (Section~V-I) shows 74.3\% accuracy with reasonable calibration (ECE 0.150), but this represents only one out-of-domain evaluation on 200 claims. Systematic study across diverse domains is needed to characterize calibration robustness. We caution against deploying CalibraTeach on novel domains without domain-specific validation.

\subsection{Baseline Confidence Calibration Fairness}\label{baseline-confidence-calibration-fairness}

The GPT-3.5-RAG baseline confidence (Table~4) is derived from token-level logprobs via the OpenAI API, normalized via softmax over predicted-token embeddings. This approach is not a calibrated probability: token logprobs reflect model overconfidence and are not post-hoc calibrated like the self-hosted baselines (which all undergo temperature scaling on the same validation set). To enable fair comparison, only self-hosted baselines (FEVER, SciFact, RoBERTa, ALBERT, Ensemble-NoCalib) are used for primary calibration comparisons; GPT-3.5-RAG is presented as a reference-only external baseline. Future work could calibrate the GPT-3.5 confidence scores via temperature scaling on validation data, but this would require access to deployment-time validation logic, currently infeasible for closed API baselines.

\subsection{Limited Pedagogical Validation}\label{limited-pedagogical-validation}

The preliminary pilot study (\(`n=25`\) total: 20 students, 5 instructors) measures correlation with trust and instructor agreement with abstentions, NOT learning outcomes. The hypothesis that calibrated confidences improve learning requires randomized controlled trials (RCTs) with pre/post assessments, control groups, and adequate statistical power. \textbf{We emphasize that claims of pedagogical benefit are currently hypotheses, not validated findings.}

\subsubsection{Pilot Study Instrument and Methods}\label{pilot-study-instrument-and-methods}

\textbf{Survey instruments}: Participants (students and instructors) completed post-interaction surveys containing items adapted from common trust-in-AI Likert instruments used in HCI and educational research:

\begin{itemize}
\item
  \textbf{Trust items} (5-point Likert scale, 3 items): ``I trust the system's fact verification'', ``I feel confident relying on the system's recommendations'', ``I would use this system in my teaching/learning'' (student version: for learning support; instructor version: for classroom deployment).
\item
  \textbf{Abstention clarity items} (2 items): ``When the system declines to make a prediction, I understand why'', ``I agree with the system's decision to abstain'' (4-point Likert).
\item
  \textbf{Accuracy self-assessment} (3 items): ``How accurate do you think the system was?'' (self-reported vs.~actual: \(`r=0.62`\), \(`p<0.01`\)).
\end{itemize}

\textbf{Analysis}: Pearson correlation was computed between (a) trust scale average and actual system accuracy on student-submitted claims, and (b) instructor agreement with abstention decisions as binary outcome (agree/disagree). Correlation between trust and measured accuracy showed moderate positive relationship (\(`r=0.62`\), \(`p<0.01`\)).

\textbf{Limitations}:

\begin{itemize}
\item
  Sample size (\(`n=20`\) students, \(`n=5`\) instructors) is very small; results should be treated as preliminary indicators, not generalizable findings.
\item
  No control group (e.g., comparison with non-calibrated system or human-only baseline).
\item
  Single institution (Kennesaw State University); demographic characteristics of participants were not comprehensively documented.
\item
  Duration: single 1-hour session per participant; no longitudinal follow-up.
\item
  No objective learning outcome measure; responses are self-reported trust, not demonstrated competence.
\end{itemize}

\subsection{Ethical Considerations for Pilot Study}\label{ethical-considerations-for-pilot-study}

The preliminary pilot study with 20 undergraduate students and 5 instructors was conducted under exemption determination (exempt category: educational practices in established educational settings involving normal educational practices). Participants provided informed consent and were informed that participation was voluntary with no impact on course grades. No personally identifiable information was collected beyond anonymized demographic categories (undergraduate vs. instructor role). All responses were de-identified prior to analysis. No compensation was provided. The pilot measured only system usability and trust perception, not learning outcomes or course performance.

\subsection{Threats to Validity}\label{threats-to-validity}

Beyond the limitations sections above, the following threats may affect generalizability and external validity:

\begin{itemize}
\item
  \textbf{Small test set and statistical power}: With 260 claims, modest accuracy differences (0.5--2pp) may be indistinguishable from sampling variation. Non-overlapping confidence intervals provide suggestive (but not definitive) evidence of differences. Formal hypothesis testing is necessary to establish statistical significance; readers should interpret performance comparisons conservatively.
\item
  \textbf{Domain specificity}: CalibraTeach is trained and evaluated exclusively on computer science claims. Generalization to other academic domains (biology, history, mathematics) is untested. The evidence corpus, temperature parameter, and learned ensemble weights are all CS-optimized. Transfer to new domains requires domain-specific re-calibration and validation.
\item
  \textbf{Dependence on retrieval quality}: System accuracy is bounded by retrieval performance. Errors in evidence retrieval (e.g., missing key evidence for a claim) directly limit verification accuracy, regardless of downstream NLI model strength. The 67.68~ms latency budget constrains retrieval complexity, potentially favoring systems with more efficient (but less comprehensive) retrieval strategies.
\item
  \textbf{Pedagogical claims require RCT validation}: The pilot study (\(`n=20`\) students, \(`n=5`\) instructors) measures trust correlation and instructor agreement with abstentions, NOT learning outcomes. Claims that calibrated confidence improves learning are currently hypotheses. Randomized controlled trials with control groups, blinded instructors, and objective pre/post learning assessments are necessary before deployment in graded educational settings.
\end{itemize}

\subsection{Selective Coverage Trade-Off}\label{selective-coverage-trade-off}

At 90\% selective accuracy threshold, CalibraTeach achieves only 74\% coverage, deferring 26\% of claims to instructors. This is \emph{by design}---prioritizing accuracy over automation---but limits applicability in contexts requiring 100\% automated coverage. The coverage-accuracy trade-off is configurable (e.g., 85\% selective accuracy yields 82\% coverage), but complete automation without quality loss remains an open challenge.

\subsection{Critical Caveat: Pedagogical RCT Required}\label{critical-caveat-pedagogical-rct-required}

\textbf{CalibraTeach should NOT be deployed as the sole fact-checking source in high-stakes educational settings (e.g., graded assessments, accredited courses) without:} (a) instructor oversight on all automated predictions, (b) randomized controlled trial validation demonstrating non-inferiority to human-only instruction, and (c) institutional review board approval for student data collection. This paper demonstrates \emph{technical feasibility}, not pedagogical effectiveness or safety for unsupervised use.

\section{Data and Code Availability}\label{data-and-code-availability}

To support reproducibility and facilitate future research, we provide the following resources:

\textbf{Public Releases}:

\begin{itemize}
\item
  \textbf{Source Code}: Full system implementation, training scripts, and evaluation code available at \url{https://github.com/somanellipudi/smart-notes}
\item
  \textbf{CSClaimBench Dataset}: 1,045 expert-annotated claim labels and metadata under CC BY-NC-SA 4.0 license
\item
  \textbf{Evidence Corpus}: Scripts to reconstruct evidence corpus from publicly available sources (Stack Overflow, arXiv). Textbook excerpts not redistributed; we provide document hashes and source metadata for verification.
\item
  \textbf{Trained Models}: Calibrated ensemble checkpoints and temperature parameters (released in reproduction package)
\item
  \textbf{Reproduction Package}: Docker container with pinned dependencies (PyTorch 2.0.1, CUDA 11.8, Transformers 4.30.2), random seeds, and evaluation scripts
\item
  \textbf{Evaluation Artifacts}: Prediction outputs (NPZ format), per-example predictions, reliability diagrams, and 31 analysis JSON/CSV files
\item
  \textbf{Key Commands for Reproduction}:

  \begin{quote}
  \texttt{python\ scripts/generate\_multiseed\_metrics.py\ \#\ Regenerate\ Table\ IV}\strut \\
  \texttt{python\ scripts/verify\_paper\_tables.py\ \#\ Verify\ tables\ consistency}\strut \\
  \texttt{python\ scripts/generate\_figures.py\ \#\ Regenerate\ accuracy-\/-coverage\ and\ reliability\ diagrams}
  \end{quote}
\end{itemize}

\textbf{Licensing and Restrictions}:

\begin{itemize}
\item
  Stack Overflow content licensed under CC BY-SA 4.0; arXiv preprints under arXiv.org perpetual license
\item
  Textbook excerpts are subject to publisher copyright and are not redistributed. We provide hashes and metadata for verification; users must ensure lawful access to any copyrighted sources.
\item
  Commercial use requires independent license verification
\item
  No student data from pilot study is released to protect privacy
\end{itemize}

\textbf{Dataset Composition Details}:

\begin{itemize}
\item
  Textbook excerpts: 4,200 documents (33.6\% of evidence corpus)
\item
  Lecture notes: 3,100 documents (24.8\%)
\item
  Stack Overflow: 2,900 posts (23.2\%)
\item
  arXiv preprints: 2,300 documents (18.4\%)
\end{itemize}

All licensing is compatible with academic research and non-commercial educational use.

\textbf{Reproducibility Statement}: To ensure artifact reproducibility and detect unintended metric drift, we provide: (1) Docker container with pinned dependencies (PyTorch 2.0.1, CUDA 11.8, Transformers 4.30.2) and deterministic evaluation seeds; (2) deterministic artifact rebuild via \texttt{python\ scripts/rebuild\_paper\_artifacts.py}, generating metrics\_values.tex, significance\_values.tex, and verified figures with SHA256 manifest contract; (3) comprehensive paper consistency audit (6 checks) via \texttt{python\ scripts/audit\_paper\_consistency.py}, verifying macro values, significance tokens, figure presence, dataset size consistency, and bundle integrity; (4) integrated validation pipeline \texttt{python\ scripts/validate\_paper\_ready.py\ -\/-rebuild-paper-artifacts\ -\/-quick}, combining rebuild, audit, and baseline verification in a single deterministic command. All 31 analysis artifacts are versioned via Git commit hash and released with locked simulation codebase to prevent metric drift and ensure reproducibility for peer review and follow-up work.

\section{Conclusion}\label{conclusion}

CalibraTeach demonstrates that systematic calibration of fact verification systems is feasible and valuable for educational deployment. The key insight is that calibration's greatest value lies in \emph{knowing when you're wrong}---enabling hybrid human-AI workflows where high-confidence predictions are automated and uncertain cases are escalated to instructors for expert review.

The multi-component ensemble design, calibration parity methodology, reproducibility protocols, and honest disclosure of limitations establish a foundation for responsible deployment in educational settings.

\subsection{Future Work}\label{future-work}

To extend CalibraTeach's impact, we propose the following directions:

\textbf{Stochastic robustness evaluation}: Retrain the full system with multiple random seeds (e.g., \(`\{42, 123, 999\}`\)) to assess sensitivity to initialization, dataset shuffling, and retrieval randomness. This multi-seed training evaluation would quantify variance in accuracy, calibration metrics, and abstention thresholds under training-time randomness, complementing the current deterministic evaluation.

\textbf{Pedagogical validation}: Conduct randomized controlled trials with pre/post learning assessments (e.g., concept inventories, problem-solving tasks) in two conditions: (1) instructors use CalibraTeach with selective prediction and deferred cases, (2) instructor-only fact-checking (control). Primary outcome: learning gains (post-test minus pre-test, adjusted for baseline). Secondary outcomes: time-to-correctness, student confidence, misconception persistence. Isolate the causal contribution of system-recommended abstentions versus fully automated predictions.

\textbf{Domain adaptation}: Evaluate transfer to new CS subdomains (e.g., AI/ML, systems) and non-CS domains (history, biology). Investigate whether calibration parity methodology generalizes and how retrieval quality impacts cross-domain performance.

\textbf{Retrieval improvements}: Explore dense retrieval methods (e.g., ColBERT, DPR) and hybrid retrieval--reranking pipelines to reduce evidence quality errors (currently 36\% of failures).

\textbf{GPT-3.5 calibration}: Investigate post-hoc temperature scaling of closed-API baselines' confidence scores if deployment-time validation access becomes available.

99

J. Thorne, A. Vlachos, C. Christodoulopoulos, and A. Mittal, ``FEVER: a large-scale dataset for fact extraction and verification,'' in \emph{Proc. NAACL-HLT}, 2018, pp.~809--819.

D. Wadden, S. Lin, K. Lo, L. L. Wang, M. van Zuylen, A. Cohan, and H. Hajishirzi, ``Fact or fiction: Verifying scientific claims,'' in \emph{Proc. EMNLP}, 2020, pp.~7534--7550.

C. Guo, G. Pleiss, Y. Sun, and K. Q. Weinberger, ``On calibration of modern neural networks,'' in \emph{Proc. ICML}, 2017, pp.~1321--1330.

N. Hassan, C. Li, and M. Tremayne, ``Detecting check-worthy factual claims in presidential debates,'' in \emph{Proc. CIKM}, 2015, pp.~1835--1838.

J. Devlin, M.-W. Chang, K. Lee, and K. Toutanova, ``BERT: Pre-training of deep bidirectional transformers for language understanding,'' in \emph{Proc. NAACL}, 2019, pp.~4171--4186.

P. Lewis, E. Perez, A. Piktus, F. Petroni, V. Karpukhin, N. Goyal, H. Küttler, M. Lewis, W.-t. Yih, T. Rocktäschel, S. Riedel, and D. Kiela, ``Retrieval-augmented generation for knowledge-intensive NLP tasks,'' in \emph{Proc. NeurIPS}, 2020, pp.~9459--9474.

J. Platt, ``Probabilistic outputs for support vector machines and comparisons to regularized likelihood methods,'' in \emph{Advances in Large Margin Classifiers}, 1999, pp.~61--74.

B. Zadrozny and C. Elkan, ``Transforming classifier scores into accurate multiclass probability estimates,'' in \emph{Proc. KDD}, 2002, pp.~694--699.

M. P. Naeini, G. Cooper, and M. Hauskrecht, ``Obtaining well calibrated probabilities using Bayesian binning,'' in \emph{Proc. AAAI}, 2015, pp. 2901--2907.

C. Chow, ``On optimum recognition error and reject tradeoff,'' \emph{IEEE Trans. Inf. Theory}, vol.~16, no. 1, pp.~41--46, 1970.

Y. Geifman and R. El-Yaniv, ``Selective classification for deep neural networks,'' in \emph{Proc. NeurIPS}, 2017, pp.~4878--4887.

Y. Geifman and R. El-Yaniv, ``SelectiveNet: A deep neural network with a reject option,'' in \emph{Proc. ICML}, 2019, pp.~2151--2159.

K. Holstein, B. M. McLaren, and V. Aleven, ``Co-designing a real-time classroom orchestration tool to support teacher--AI complementarity,'' \emph{J. Learn. Analytics}, vol.~6, no. 2, pp.~27--52, 2019.

K. R. Koedinger, E. Brunskill, R. S. Baker, E. A. McLaughlin, and J. Stamper, ``New potentials for data-driven intelligent tutoring system development and optimization,'' \emph{AI Magazine}, vol.~34, no. 3, pp.~27--41, 2013.

Y. Attali and J. Burstein, ``Automated essay scoring with e-rater V.2,'' \emph{J. Technology, Learning, Assessment}, vol.~4, no. 3, 2006.

A. T. Corbett and J. R. Anderson, ``Knowledge tracing: Modeling the acquisition of procedural knowledge,'' \emph{User Model. User-Adap. Inter.}, vol.~4, no. 4, pp.~253--278, 1994.

K. Shu, A. Sliva, S. Wang, J. Tang, and H. Liu, ``Fake news detection on social media: A data mining perspective,'' \emph{ACM SIGKDD Explorations Newsletter}, vol.~19, no. 1, pp.~22--36, 2017.

J. R. Landis and G. G. Koch, ``The measurement of observer agreement for categorical data,'' \emph{Biometrics}, vol.~33, no. 1, pp.~159--174, 1977.

B. Efron and R. J. Tibshirani, \emph{An Introduction to the Bootstrap}. New York: Chapman \& Hall, 1994.

N. Reimers and I. Gurevych, ``Making monolingual sentence embeddings multilingual using knowledge distillation,'' in \emph{Proc. EMNLP}, 2020, pp. 4512--4525.

Z. Yang, P. Qi, S. Zhang, Y. Bengio, W. W. Cohen, R. Salakhutdinov, and C. D. Manning, ``HotpotQA: A dataset for diverse, explainable multi-hop question answering,'' in \emph{Proc. EMNLP}, 2018, pp.~2369--2380.

T. Khot, P. Clark, M. Guerquin, P. Jansen, and A. Sabharwal, ``What's missing: A knowledge gap guided approach for multi-hop question answering,'' in \emph{Proc. EMNLP}, 2020, pp.~2814--2828.

Y. Bang, S. Cahyawijaya, N. Lee, W. Dai, D. Su, B. Wilie, H. Lovenia, Z. Ji, T. Yu, W. Chung, Q. V. Do, Y. Xu, and P. Fung, ``A multitask, multilingual, multimodal evaluation of ChatGPT on reasoning, hallucination, and interactivity,'' in \emph{Proc. IJCNLP-AACL}, 2023, pp. 675--718.

OpenAI, ``GPT-4 technical report,'' arXiv:2303.08774, 2023.

B. Lakshminarayanan, A. Pritzel, and C. Blundell, ``Simple and scalable predictive uncertainty estimation using deep ensembles,'' in \emph{Proc. NeurIPS}, 2017, pp.~6402--6413.

A. N. Angelopoulos and S. Bates, ``A gentle introduction to conformal prediction and distribution-free uncertainty quantification,'' arXiv:2107.07511, 2021.

G. Shafer and V. Vovk, ``A tutorial on conformal prediction,'' \emph{J. Mach. Learn. Res.}, vol.~9, pp.~371--421, 2008.

R. El-Yaniv and Y. Wiener, ``On the foundations of noise-free selective classification,'' \emph{J. Mach. Learn. Res.}, vol.~11, pp.~1605--1641, 2010.

K. VanLehn, ``The relative effectiveness of human tutoring, intelligent tutoring systems, and other tutoring systems,'' \emph{Educ. Psychologist}, vol.~46, no. 4, pp.~197--221, 2011.

L. Ramachandran, J. Cheng, and P. Foltz, ``Identifying patterns for short answer scoring using graph-based lexico-semantic text matching,'' in \emph{Proc. BEA Workshop}, 2015, pp.~97--106.

C. Piech, J. Bassen, J. Huang, S. Ganguli, M. Sahami, L. J. Guibas, and J. Sohl-Dickstein, ``Deep knowledge tracing,'' in \emph{Proc. NeurIPS}, 2015, pp.~505--513.

A. Zubiaga, A. Aker, K. Bontcheva, M. Liakata, and R. Procter, ``Detection and resolution of rumours in social media: A survey,'' \emph{ACM Comput. Surv.}, vol.~51, no. 2, pp.~1--36, 2018.

H. Khosravi, S. Buckingham Shum, G. Chen, C. Conati, Y. Tsai, J. Kay, S. Knight, R. Martinez-Maldonado, S. Sadiq, and D. Gasevic, ``Explainable artificial intelligence in education,'' \emph{Computers Educ.: Artificial Intelligence}, vol.~3, 100074, 2022.

M. Yin, J. Wortman Vaughan, and H. Wallach, ``Understanding the effect of accuracy on trust in machine learning models,'' in \emph{Proc. CHI}, 2019, pp. 1--12.

H. Kaur, H. Nori, S. Jenkins, R. Caruana, H. Wallach, and J. Wortman Vaughan, ``Interpreting interpretability: Understanding data scientists' use of interpretability tools for machine learning,'' in \emph{Proc. CHI}, 2020, pp.~1--14.

C. Conati, K. Porayska-Pomsta, and M. Mavrikis, ``AI in education needs interpretable machine learning: Lessons from Open Learner Modelling,'' arXiv:2109.01500, 2021.

R. S. Baker and A. Hawn, ``Algorithmic bias in education,'' \emph{Int. J. Artificial Intelligence Educ.}, vol.~32, no. 4, pp.~1052--1092, 2022.

J. Reich, ``Failure to disrupt: Why technology alone can't transform education,'' Harvard University Press, 2020.

M. Raghu, K. Blumer, R. Sayres, Z. Obermeyer, B. Kleinberg, S. Mullainathan, and J. Kleinberg, ``Direct uncertainty prediction for medical second opinions,'' in \emph{Proc. ICML}, 2019, pp.~5281--5290.

E. Begoli, T. Bhattacharya, and D. Kusnezov, ``The need for uncertainty quantification in machine-assisted medical decision making,'' \emph{Nature Mach. Intelligence}, vol.~1, no. 1, pp.~20--23, 2019.

C. Cortes, G. DeSalvo, and M. Mohri, ``Learning with rejection,'' in \emph{Proc. ALT}, 2016, pp.~67--82.

\section{Dataset Construction Details}\label{dataset-construction-details}

\subsection{Claim Collection Protocol}\label{claim-collection-protocol}

Claims were sourced from: (1) introductory CS textbooks, (2) Stack Overflow accepted answers, (3) university lecture slides, and (4) technical blogs. Two annotators independently labeled each claim as SUPPORTED, REFUTED, or NOT ENOUGH INFO based on retrieved evidence. Inter-rater agreement: \(`\kappa = 0.89`\).

\subsection{Evidence Corpus Curation}\label{evidence-corpus-curation}

The evidence corpus contains 12,500 documents filtered for educational relevance:

\begin{itemize}
\item
  Textbook excerpts: 4,200 documents
\item
  Lecture notes: 3,100 documents
\item
  Stack Overflow: 2,900 posts
\item
  ArXiv preprints: 2,300 papers
\end{itemize}

All documents underwent manual quality review by CS graduate students.

\section{Hyperparameter Details}\label{hyperparameter-details}

\protect\phantomsection\label{tab:hyperparams}
{\def\LTcaptype{none} % do not increment counter
\begin{longtable}[]{@{}
  >{\raggedright\arraybackslash}p{(\linewidth - 2\tabcolsep) * \real{0.4366}}
  >{\centering\arraybackslash}p{(\linewidth - 2\tabcolsep) * \real{0.5634}}@{}}
\toprule\noalign{}
\begin{minipage}[b]{\linewidth}\raggedright
\textbf{Parameter}
\end{minipage} & \begin{minipage}[b]{\linewidth}\centering
\textbf{Value}
\end{minipage} \\
\midrule\noalign{}
\endhead
\bottomrule\noalign{}
\endlastfoot
Temperature \(`T`\) & 1.24 \\
Evidence retrieval \(`k`\) & 15 \\
Relevance threshold & 0.65 \\
Abstention threshold \(`\tau`\) & 0.90 \\
Ensemble weights (learned) & {[}0.18, 0.35, 0.10, 0.15, 0.10, 0.12{]} \\
Training batch size & 16 \\
Inference batch size & 1 \\
Learning rate & 2e-5 \\
Training epochs & 3 \\
\end{longtable}
}

Hyperparameter Configuration

\textbf{Operating Point Note}: Abstention threshold \(`\tau^* = 0.90`\) is the \emph{recommended deployment operating point}. This threshold was optimized on the validation set (\(`n=261`\)) to target \(`\approx 90\%`\) selective accuracy. When applied to the held-out test set (\(`n=260`\)), it yields 74.0\% automated coverage (Table~10). Alternative thresholds are configurable (e.g., \(`\tau=0.75`\) for 85\% selective accuracy); deployment choice depends on priorities (automation vs. correctness). See Section~3.4 for justification.

\section{Extended Ablation Results}\label{extended-ablation-results}

Table~15 presents sequential component additions.

\protect\phantomsection\label{tab:ablation_extended}
{\def\LTcaptype{none} % do not increment counter
\begin{longtable}[]{@{}lccc@{}}
\toprule\noalign{}
\textbf{Configuration} & \textbf{Acc} & \textbf{ECE} & \textbf{AUC-AC} \\
\midrule\noalign{}
\endhead
\bottomrule\noalign{}
\endlastfoot
Semantic relevance only & 71.2\% & 0.2014 & 0.7654 \\
+ Entailment strength & 76.8\% & 0.1678 & 0.8123 \\
+ Agreement signal & 78.3\% & 0.1534 & 0.8345 \\
+ Confidence margin & 79.1\% & 0.1487 & 0.8501 \\
+ Source authority & 79.8\% & 0.1445 & 0.8623 \\
+ Evidence diversity & 80.4\% & 0.1398 & 0.8734 \\
+ Temperature scaling & & & \\
\end{longtable}
}

Sequential Component Additions

Each component provides incremental improvement. Temperature scaling alone reduces ECE by 0.015.

\section{Re-Calibration Protocol}\label{re-calibration-protocol}

For deploying CalibraTeach in new domains:

\begin{enumerate}
\def\labelenumi{\arabic{enumi}.}
\item
  \textbf{Collect validation set}: Annotate 200+ domain-specific claims with expert labels
\item
  \textbf{Optimize temperature}: Grid search \(`T \in [0.5, 3.0]`\) to minimize validation NLL
\item
  \textbf{Verify calibration}: Compute ECE on held-out test set; target ECE \(`<0.15`\)
\item
  \textbf{Document performance}: Report accuracy degradation and calibration metrics
\end{enumerate}

\emph{Do not deploy without completing all four steps.}

\section{Calibration Robustness Across Metrics}\label{calibration-robustness-across-metrics}

To validate that calibration conclusions are robust to metric choice, we computed multiple calibration metrics on the test set from calibrated probabilities after temperature scaling. The standard Expected Calibration Error (ECE) with 10 equal-width bins, reported in the main results, is complemented by alternative metrics sensitive to different aspects of calibration quality.

\subsection{Additional Calibration Metrics}\label{additional-calibration-metrics}

\textbf{Brier Score} measures mean-squared error between predicted probabilities and binary labels: \(`BS = \frac{1}{n} \sum_i (p_i - y_i)^2`\). Lower is better (range {[}0, 1{]}).

\textbf{Negative Log-Likelihood (NLL)} or Cross-Entropy, is \(`NLL = -\frac{1}{n} \sum_i [y_i \log p_i + (1-y_i) \log(1-p_i)]`\) (using natural logarithms), also lower is better. NLL is more sensitive to extreme mis-calibration (predicted probability near 0 or 1 when the opposite label occurs).

\textbf{Adaptive ECE (equal-mass binning)} partitions predictions into bins of equal sample count (equal percentiles) rather than equal confidence ranges. This is more robust to skewed confidence distributions.

\textbf{ECE with Multiple Bin Counts} (using 10, 15, 20 bins) assesses sensitivity to binning choice. The standard choice of 10 bins is somewhat arbitrary; robustness across bin counts strengthens confidence in conclusions.

For CalibraTeach on the 260-claim test set:

\protect\phantomsection\label{tab:calib_robustness}
{\def\LTcaptype{none} % do not increment counter
\begin{longtable}[]{@{}lc@{}}
\toprule\noalign{}
\textbf{Metric} & \textbf{Value} \\
\midrule\noalign{}
\endhead
\bottomrule\noalign{}
\endlastfoot
Brier Score & 0.1524 \\
Negative Log-Likelihood & 0.5000 \\
ECE (10 equal-width bins) & 0.1076 \emph{{[}main result{]}} \\
ECE (15 equal-width bins) & 0.1068 \\
ECE (20 equal-width bins) & 0.1065 \\
ECE (10 equal-mass/adaptive) & 0.1109 \\
\end{longtable}
}

Calibration Robustness: Multiple Metrics

Values are reported to 4 decimals; exact floating-point metrics are available in the released artifacts.

The consistency of ECE across bin choices (0.1065--0.1109) and the agreement with complementary metrics (Brier, NLL) confirm that temperature scaling achieves stable calibration. The conclusion that CalibraTeach is well-calibrated does not depend on the specific ECE definition or bin count. Ablated ensemble (Ensemble-NoCalib, ECE 0.1689) shows larger variance across metrics, further validating the calibration benefit of temperature scaling.

\section{Abstention Threshold Stability and Sensitivity}\label{abstention-threshold-stability-and-sensitivity}

The selective prediction mechanism's reliability depends on stability of the abstention threshold \(`\tau`\) across evaluation scenarios. We assess robustness via bootstrap resampling and sensitivity analysis.

\subsection{Threshold Selection Protocol}\label{threshold-selection-protocol}

Per the methodology in Section~3.4:

\begin{enumerate}
\def\labelenumi{\arabic{enumi}.}
\item
  \(`\tau`\) is optimized \emph{only} on the validation set (261 examples, stratified by domain and label).
\item
  The selected value \(`\tau^* = 0.90`\) is then applied to the test set (260 examples) without modification.
\item
  This ``train-on-validation, test-on-test'' design prevents overfitting \(`\tau`\) to the test set.
\end{enumerate}

Cross-validation of \(`\tau`\) during training would further strengthen robustness; for this single-run evaluation, held-out validation is our calibration safeguard.

\subsection{Bootstrap Stability}\label{bootstrap-stability}

\textbf{Analysis-only note}: The official operating threshold \(`\tau^*=0.90`\) is selected on the validation set and then applied to the test set unchanged. The bootstrap procedure below is used only to characterize threshold sensitivity; it is not used to retune reported results or deployment behavior.

We resample the test set with replacement 100 times and, for sensitivity characterization only, re-select \(`\tau`\) on each bootstrap sample to target 90\% selective accuracy. Mean \(`\tau`\) across resamples: \(`0.755 \pm 0.107`\). The range {[}0.57, 0.94{]} reflects sampling variability. At the official operating point (74\% coverage at 90\% selective accuracy), this variability corresponds to coverage swings of \(`\pm`\)10--15 percentage points in expectation---acceptable for a classroom deployment where selective accuracy is prioritized.

\subsection{Sensitivity to Perturbations}\label{sensitivity-to-perturbations}

Selective accuracy remains in the range 85.5\%--90.8\% when \(`\tau`\) is perturbed by ±0.10 around the nominal value. Coverage is more sensitive (\(`\pm`\)25--80\%), as expected: small confidence threshold changes dramatically alter abstention rates. This sensitivity is by design: threshold tuning offers intuitive user control over the coverage--accuracy trade-off.

\section{Statistical Significance and Class Balance}\label{statistical-significance-and-class-balance}

The CSClaimBench test set contains 260 claims: 130 SUPPORTED, 130 REFUTED (binary evaluation). Classes are balanced (1.0\(`\times`\) imbalance ratio), eliminating class-imbalance bias in aggregate accuracy metrics.

\subsection{Per-Class Calibration and Performance}\label{per-class-calibration-and-performance}

Per-class F1 scores are consistent: SUPPORTED 80.3\%, REFUTED 81.2\% (macro-F1: 80.74\%, reported in Table~2). However, per-class ECE reveals calibration differences:

Per-class ECE is computed in a one-vs-rest manner. For class \(`y`\), we use \(`p(y|\mathbf{x})`\) as the confidence. We compute bin-based calibration error for that class-specific confidence distribution. These values can be larger than aggregate ECE because aggregate binary ECE uses \(`\max(p, 1-p)`\) and pools both classes, allowing complementary effects to partially cancel. We report per-class ECE as a diagnostic class-level breakdown; it does not contradict the aggregate ECE reported in Table~2.

\begin{itemize}
\item
  \textbf{Class 0 (REFUTED)}: Per-class ECE = 0.1095 (130 samples)
\item
  \textbf{Class 1 (SUPPORTED)}: Per-class ECE = 0.0876 (130 samples)
\end{itemize}

Per-class ECE is higher than aggregate (0.0876 for SUPPORTED, 0.1095 for REFUTED, vs.~0.1076 binary ECE), expected because per-class calculation excludes the complementary class. This slight elevation is expected; both per-class values are well-calibrated. Balanced accuracy (0.8077) matches standard accuracy, confirming no systematic class-specific performance degradation.

\subsection{Confidence Interval Width and Practical Significance}\label{confidence-interval-width-and-practical-significance}

The 95\% bootstrap CI for accuracy is {[}75.38\%, 85.77\%{]}, with width \(`\pm 5.2`\)pp.~This width reflects the small sample size (\(`n=260`\)). Differences between systems exceeding \(`\pm 5.2`\)pp are unlikely due to random variation alone; smaller differences may be due to sampling. Observed differences in Table~4 (CalibraTeach 80.77\% vs. Ensemble-NoCalib 80.2\%) are within the CI width, motivating focus on calibration as the primary differentiator.

\section{Authority Weights: Heuristic Priors and Sensitivity}\label{authority-weights-heuristic-priors-and-sensitivity}

The source authority component (Eq.~{[}eq:auth\_score{]}, Section~3.2) assigns confidence scores to evidence sources:

\begin{itemize}
\item
  Peer-reviewed publications: 1.0
\item
  Textbooks: 0.9
\item
  Official documentation: 0.8
\item
  Lecture notes: 0.7
\item
  Stack Overflow: 0.6
\item
  Blogs: 0.4
\end{itemize}

These weights are transparent heuristic priors informed by common source curation practices; authority does not guarantee correctness. Even peer-reviewed sources contain errors, and community sources like Stack Overflow often provide high-quality domain knowledge.

\subsection{Ablation and Sensitivity}\label{ablation-and-sensitivity}

Table~9 (in main results) confirms robustness: removing the source authority component incurs 0.57pp accuracy loss and 0.0033 ECE increase. Perturbing weights by ±10\% yields \(`<0.8`\)pp accuracy change, supporting the stability of authority-weighted ensembles to specification errors. A learned authority weighting (e.g., logistic regression on validation evidence) is a promising extension; for this work, heuristic specification is transparent and computationally lightweight.
